% ray-tracing-cellular-computing.tex
% Ray-Tracing as Cellular Computation: Simultaneous Optical,
% Chromatographic, and Circuit Observation Through Volumetric
% Partition Traversal
% Author: Kundai Farai Sachikonye
% Date: 2026

\documentclass[11pt,twocolumn,a4paper]{article}

% ─── Packages ────────────────────────────────────────────────────────
\usepackage[utf8]{inputenc}
\usepackage[T1]{fontenc}
\usepackage{lmodern}
\usepackage{amsmath,amssymb,amsthm,mathtools}
\usepackage{physics}
\usepackage{bm}
\usepackage{graphicx}
\usepackage{xcolor}
\usepackage{booktabs}
\usepackage{array}
\usepackage{multirow}
\usepackage{enumitem}
\usepackage{float}
\usepackage{caption}
\usepackage{subcaption}
\usepackage[colorlinks=true,linkcolor=blue!60!black,citecolor=blue!60!black,urlcolor=blue!60!black]{hyperref}
\usepackage[margin=1.8cm,top=2.2cm,bottom=2.2cm]{geometry}
\usepackage{natbib}
\usepackage{fancyhdr}
\usepackage{titlesec}
\usepackage{microtype}
\usepackage{algorithm}
\usepackage{algpseudocode}
\usepackage{listings}
\usepackage{tabularx}
\usepackage{tcolorbox}

% ─── Listings Style ─────────────────────────────────────────────────
\lstdefinestyle{glsl}{
  language=C,
  basicstyle=\ttfamily\scriptsize,
  keywordstyle=\color{blue!70!black}\bfseries,
  commentstyle=\color{green!50!black}\itshape,
  stringstyle=\color{red!60!black},
  numbers=left,
  numberstyle=\tiny\color{gray},
  numbersep=4pt,
  frame=single,
  breaklines=true,
  columns=fullflexible,
  morekeywords={vec2,vec3,vec4,mat2,mat3,mat4,sampler2D,sampler3D,
    uniform,varying,in,out,layout,location,precision,
    highp,mediump,lowp,texture,gl_FragColor,void,float,int,
    ivec3,uvec3,bool}
}

\lstdefinestyle{pseudo}{
  basicstyle=\ttfamily\scriptsize,
  keywordstyle=\color{blue!70!black}\bfseries,
  commentstyle=\color{green!50!black}\itshape,
  frame=single,
  breaklines=true,
  columns=fullflexible
}

% ─── Theorem Environments ────────────────────────────────────────────
\newtheorem{theorem}{Theorem}[section]
\newtheorem{lemma}[theorem]{Lemma}
\newtheorem{corollary}[theorem]{Corollary}
\newtheorem{proposition}[theorem]{Proposition}
\newtheorem{definition}[theorem]{Definition}
\newtheorem{remark}[theorem]{Remark}
\newtheorem{axiom}{Axiom}
\newtheorem{example}[theorem]{Example}

% ─── Custom Commands ─────────────────────────────────────────────────
\newcommand{\kb}{k_{\mathrm{B}}}
\newcommand{\Ocat}{\hat{O}_{\mathrm{cat}}}
\newcommand{\Ophys}{\hat{O}_{\mathrm{phys}}}
\newcommand{\Sk}{S_k}
\newcommand{\St}{S_t}
\newcommand{\Se}{S_e}
\newcommand{\BigO}{\mathcal{O}}
\newcommand{\R}{\mathbb{R}}
\newcommand{\N}{\mathbb{N}}
\newcommand{\Z}{\mathbb{Z}}
\newcommand{\CC}{\mathbb{C}}
\newcommand{\Hilb}{\mathcal{H}}
\newcommand{\Part}{\mathcal{P}}
\newcommand{\Cat}{\mathcal{C}}
\newcommand{\Tex}{\mathcal{T}}
\newcommand{\Obs}{\mathcal{O}}
\newcommand{\Rend}{\mathcal{R}}
\newcommand{\Meas}{\mathcal{M}}
\newcommand{\taup}{\tau_{\mathrm{p}}}
\newcommand{\tauem}{\tau_{\mathrm{em}}}
\newcommand{\dcat}{d_{\mathrm{cat}}}
\newcommand{\dS}{d_{\mathrm{S}}}
\newcommand{\etacell}{\eta_{\mathrm{cell}}}
\newcommand{\Vcell}{V_{\mathrm{cell}}}
\newcommand{\muabs}{\mu_{\mathrm{a}}}
\newcommand{\Jray}{J_{\mathrm{ray}}}
\newcommand{\cS}{c_{\mathrm{S}}}
\DeclareMathOperator{\Tr}{Tr}
\DeclareMathOperator{\diag}{diag}
\DeclareMathOperator{\sgn}{sgn}
\DeclareMathOperator{\Vis}{Vis}
\DeclareMathOperator{\Ent}{H}

% ─── Page Style ──────────────────────────────────────────────────────
\pagestyle{fancy}
\fancyhf{}
\fancyhead[L]{\small\textit{Ray-Tracing as Cellular Computation}}
\fancyhead[R]{\small\thepage}
\renewcommand{\headrulewidth}{0.4pt}

% ─── Title Section Formatting ────────────────────────────────────────
\titleformat{\section}{\large\bfseries}{\thesection.}{0.5em}{}
\titleformat{\subsection}{\normalsize\bfseries}{\thesubsection.}{0.5em}{}
\titleformat{\subsubsection}{\normalsize\itshape}{\thesubsubsection.}{0.5em}{}

% ═════════════════════════════════════════════════════════════════════
\begin{document}

\twocolumn[
\begin{@twocolumnfalse}
\begin{center}

{\LARGE\bfseries Ray-Tracing as Cellular Computation:\\[4pt]
Simultaneous Optical, Chromatographic, and Circuit\\[4pt]
Observation Through Volumetric Partition Traversal}

\vspace{12pt}

{\large Kundai Farai Sachikonye}

\vspace{6pt}

{\small\itshape Independent Researcher \\[0.3em]
\texttt{kundai.sachikonye@bitspark.com}}

\vspace{6pt}

{\small April 2026}

\vspace{14pt}

\begin{minipage}{0.92\textwidth}
\textbf{Abstract.}
We prove from first principles that a single ray marching through
cellular cytoplasm simultaneously computes three physically distinct
observations at every integration step: optical absorption (light
transport via Beer--Lambert attenuation), chromatographic retention
(molecular interaction via categorical distance), and circuit current
flow (biochemical reaction network via Kirchhoff's laws). We establish
the \emph{Triple Observation Identity}: these three computations are
not independent processes performed in parallel but a single partition
observation expressed in three equivalent representations, unified by a
theorem we call the \emph{Triple Equivalence}. The cytoplasm, with its
heterogeneous molecular composition and spatial organization, is shown
to function as a three-dimensional chromatographic column whose local
retention properties are encoded by partition coordinates
$(n,\ell,m,s)$ at each spatial position. The intracellular reaction
network is shown to be isomorphic to an electrical circuit whose node
potentials are chemical potentials and whose conductances are reaction
rate constants. The ray tracer computes all three representations
simultaneously because they share a common mathematical structure: the
absorption coefficient, the inverse retention distance, and the
conductance are proportional, with proportionality constants determined
by the Triple Equivalence.

We further demonstrate that the eight oscillator classes spanning
cellular dynamics --- from protein folding vibrations at $10^{13}$\,Hz
to circadian rhythm at $10^{-5}$\,Hz --- each manifest as a distinct
type of ray--matter interaction (absorption, scattering, gated
transmission, reflection, amplification, state switching, oscillatory
modulation, and parametric drift). The accumulated phase shift along a
ray path encodes the complete reaction history of every oscillator the
ray encounters. Interference between multiple rays launched through the
cellular volume yields a visibility parameter $\Vcell$ that is
identical to the cellular coherence index $\etacell$: high visibility
(constructive interference from synchronized oscillators) indicates
healthy cells; low visibility (decoherent oscillators) indicates
pathology. We derive the holographic reconstruction of
three-dimensional cellular structure from a two-dimensional observation
via angular spectrum back-propagation, establish that the
reconstructed volumetric field encodes the complete partition state at
every three-dimensional position, and prove that a GPU ray tracer
running in a web browser simultaneously reconstructs morphology,
metabolic flow maps, reaction network state, and health diagnostics
from a single volumetric traversal. We extend the framework to
cytoplasmic fluid dynamics, showing that Kirchhoff's laws in the
continuous limit yield the Stokes equations appropriate for the
cellular interior ($Re \sim 10^{-4}$), so the ray tracer implicitly
solves for cytoplasmic flow. The complete five-pass GPU pipeline ---
dual-pixel analysis, holographic back-propagation, triple-observation
ray march, multi-ray interference, and scalar readback --- is specified
with GLSL pseudocode and memory analysis demonstrating feasibility on
integrated GPUs with $\sim\!32$\,MB working set. The theory is
validated through five independent experimental designs testing triple
observation consistency, holographic reconstruction fidelity, coherence
diagnostic accuracy, flow field recovery, and real-time throughput.

\vspace{8pt}
\noindent\textbf{Keywords:} ray tracing, cellular computation, triple
observation identity, chromatographic transport, circuit model, volume
rendering, partition coordinates, oscillator classes, cellular
coherence, holographic reconstruction, GPU computation, digital twin
\end{minipage}

\vspace{18pt}
\end{center}
\end{@twocolumnfalse}
]

% ═════════════════════════════════════════════════════════════════════
%                    PART I: THEORETICAL FOUNDATIONS
% ═════════════════════════════════════════════════════════════════════

\section{Introduction}
\label{sec:introduction}

The living cell is simultaneously an optical medium, a chromatographic
matrix, and a biochemical circuit. Light passing through cytoplasm is
absorbed and scattered by organelles and macromolecules
\citep{bornwolf1999}. Molecules diffusing through the intracellular
space interact with stationary structures --- membranes, cytoskeletal
filaments, protein complexes --- in a manner formally analogous to
analytes traversing a chromatographic column
\citep{martinsynge1941,vandeemter1956}. The network of enzymatic
reactions, ion channel conductances, and metabolic fluxes constitutes
an electrical circuit in the thermodynamic sense, with chemical
potentials as voltages and reaction rates as currents
\citep{kirchhoff1847,onsager1931}.

These three descriptions have been developed independently across
optics, analytical chemistry, and systems biology. The central
contribution of this paper is the proof that they are not independent
descriptions but \emph{the same description} expressed in three
equivalent mathematical representations. Specifically, we prove:

\begin{enumerate}[leftmargin=*,itemsep=2pt]
\item A single ray marching through a volumetric representation of the
  cell computes optical absorption, chromatographic retention, and
  circuit current \emph{simultaneously} at every integration step.

\item The three quantities are related by proportionality constants
  derivable from first principles: $\muabs(r) \propto 1/(\tau \cdot
  \dS) \propto G(r) \cdot RT$.

\item The eight oscillator classes spanning 18 decades of cellular
  frequency (from protein folding at $10^{13}$\,Hz to circadian
  rhythm at $10^{-5}$\,Hz) each manifest as a specific type of
  ray--matter interaction.

\item Interference between rays yields a visibility parameter
  identical to the cellular coherence index (a scalar health
  diagnostic).

\item The complete computation runs in real time on an integrated GPU
  in a web browser.
\end{enumerate}

% ── Figure 1 proposal ──────────────────────────────────────────────
% FIGURE 1: Central schematic. A single ray traverses the cellular
% volume (shown as a 3D box with organelles). At each step, three
% outputs branch: optical (Beer-Lambert), chromatographic (retention),
% circuit (current). All three arrows converge to a single node
% labeled "Partition State Sigma(r)". Below: the eight oscillator
% frequency bands as colored bars.

The paper is organized as follows. Part~I
(\S\ref{sec:axioms}--\S\ref{sec:commutation}) derives the theoretical
foundations from first principles: bounded phase space axioms,
oscillatory necessity, partition coordinates, the Triple Equivalence,
the fundamental identity, S-entropy coordinates, and operator
commutation. Part~II
(\S\ref{sec:optical-transport}--\S\ref{sec:triple-identity}) establishes
the Triple Observation Identity: the main new contribution. Part~III
(\S\ref{sec:oscillator-classes}--\S\ref{sec:coherence}) develops the
eight oscillator classes as ray interaction types and derives the
coherence diagnostic from multi-ray interference. Part~IV
(\S\ref{sec:dual-pixel}--\S\ref{sec:volumetric-state}) derives
holographic three-dimensional reconstruction. Part~V
(\S\ref{sec:volumetric-data}--\S\ref{sec:pipeline}) specifies the GPU
ray tracer architecture. Part~VI
(\S\ref{sec:cytoplasmic-flow}--\S\ref{sec:viscosity}) extends the
framework to fluid dynamics. Part~VII
(\S\ref{sec:validation}--\S\ref{sec:conclusion}) presents validation,
discussion, and conclusions.

\subsection{Methodological Approach}

The argument proceeds by deriving all required structures from two
axioms (bounded phase space and categorical observation). No external
results are assumed beyond standard mathematics and physics. Every
theorem is proved within this paper. The approach is constructive:
we build the partition coordinate system, prove the equivalences,
derive the ray transport equations, and specify the GPU implementation
with sufficient detail for independent reproduction.


% ═════════════════════════════════════════════════════════════════════
\section{Axioms of Bounded Categorical Observation}
\label{sec:axioms}

We begin with two axioms from which all subsequent results derive.

\begin{axiom}[Bounded Phase Space]
\label{ax:bounded}
Every persistent dynamical system $X$ occupies a bounded region
$\Gamma \subset \R^{2f}$ of phase space (where $f$ is the number of
degrees of freedom) with finite Liouville measure:
\begin{equation}
\mu(\Gamma) = \int_{\Gamma} \prod_{i=1}^{f}
  \dd{q_i}\,\dd{p_i} < \infty.
\label{eq:liouville-finite}
\end{equation}
\end{axiom}

\noindent
\textit{Justification.} Physical systems are confined by potential
wells, conservation laws, or finite energy. Biological systems
occupy finite spatial extent and operate within bounded temperature
and concentration ranges. Computational representations use
finite-precision arithmetic. In every domain of interest, the state
space is bounded. Systems that access unbounded phase space are either
transient (they escape to infinity) or unphysical
\citep{goldstein2002,arnoldavez1968}.

\begin{axiom}[Categorical Observation]
\label{ax:categorical}
An observer with finite resolution $\epsilon > 0$ partitions the
bounded phase space $\Gamma$ into a finite number of distinguishable
\emph{categories}: disjoint cells $\{C_\alpha\}_{\alpha=1}^{K}$ with
$\bigcup_\alpha C_\alpha = \Gamma$ and $\mu(C_\alpha) \geq
\epsilon^{2f}$ for all $\alpha$. Two states $x, y \in \Gamma$ are
\emph{categorically equivalent} if and only if they belong to the same
cell: $x \sim y \iff \exists\,\alpha : x,y \in C_\alpha$.
\end{axiom}

\noindent
\textit{Justification.} Every physical measurement has finite
precision \citep{heisenberg1927}. Every digital sensor has finite bit
depth. Every classification system uses a finite number of categories.
The partition induced by finite resolution is a universal feature of
observation \citep{vonneumann1932}.

\begin{proposition}[Finite Category Count]
\label{prop:finite-categories}
Under Axioms~\ref{ax:bounded} and \ref{ax:categorical}, the number
of distinguishable categories is bounded:
\begin{equation}
K \leq \frac{\mu(\Gamma)}{\epsilon^{2f}}.
\label{eq:finite-K}
\end{equation}
\end{proposition}

\begin{proof}
Since the cells $\{C_\alpha\}$ are disjoint and each has measure at
least $\epsilon^{2f}$, we have $\mu(\Gamma) = \sum_\alpha
\mu(C_\alpha) \geq K \cdot \epsilon^{2f}$. Since $\mu(\Gamma) < \infty$
by Axiom~\ref{ax:bounded} and $\epsilon > 0$ by
Axiom~\ref{ax:categorical}, $K$ is finite.
\end{proof}

\begin{remark}
The two axioms --- bounded phase space and finite-resolution
observation --- are the only assumptions. Everything that follows
is a consequence.
\end{remark}


% ═════════════════════════════════════════════════════════════════════
\section{The Oscillatory Necessity Theorem}
\label{sec:oscillatory}

We now prove that oscillatory dynamics is the unique valid
long-term behavior for bounded, observed systems.

\begin{definition}[Dynamical Modes]
\label{def:modes}
Let $\Phi_t : \Gamma \to \Gamma$ be the time-evolution map. We
classify long-term behaviors:
\begin{enumerate}[label=(\roman*),leftmargin=*]
\item \textbf{Static:} $\exists\,x^* : \Phi_t(x^*) = x^*$ for all
  $t > 0$, and $\Phi_t(x) \to x^*$ for $\mu$-a.e.\ $x$.
\item \textbf{Monotonic:} $\exists$ observable $f : \Gamma \to \R$
  such that $f(\Phi_t(x))$ is strictly monotone in $t$ for
  $\mu$-a.e.\ $x$.
\item \textbf{Chaotic (non-recurrent):} positive Lyapunov exponent
  $\lambda_{\max} > 0$ and no recurrence: for $\mu$-a.e.\ $x$ and
  every $\epsilon > 0$, $\exists\,T$ such that
  $\|\Phi_t(x) - x\| > \epsilon$ for all $t > T$.
\item \textbf{Oscillatory:} for $\mu$-a.e.\ $x$ and every
  $\epsilon > 0$, the set
  $\{t > 0 : \|\Phi_t(x) - x\| < \epsilon\}$ is unbounded, and the
  system visits a positive-measure subset of $\Gamma$ infinitely often.
\end{enumerate}
\end{definition}

\begin{theorem}[Oscillatory Necessity]
\label{thm:oscillatory-necessity}
Under Axioms~\ref{ax:bounded} and \ref{ax:categorical}, every
persistent dynamical system exhibits oscillatory dynamics. Static,
monotonic, and non-recurrent chaotic modes are excluded.
\end{theorem}

\begin{proof}
We eliminate the three non-oscillatory modes:

\textit{Elimination of static mode.}
A system at a fixed point $x^*$ occupies a single categorical cell
$C_\alpha$. The categorical entropy is $H_{\mathrm{cat}} = 0$,
representing zero information capacity. But any system that persists
long enough to be observed must have carried information to arrive at
$x^*$, requiring $H_{\mathrm{cat}} > 0$ at some prior time. If the
system was always at $x^*$, it carries no dynamical information and is
observationally trivial --- not a persistent dynamical system. Static
equilibria are absorbing states, not persistent dynamics
\citep{strogatz2015}.

\textit{Elimination of monotonic mode.}
If $f(\Phi_t(x))$ is strictly monotone, then $f$ must be strictly
increasing or strictly decreasing along every trajectory. Since
$\Gamma$ is bounded, $f$ is bounded on $\Gamma$: $f_{\min} \leq f(x)
\leq f_{\max}$ for all $x \in \Gamma$. A strictly monotone bounded
sequence must converge. But convergence to a limit contradicts strict
monotonicity (eventually $f$ must become constant or reverse). Hence no
observable can be strictly monotone on a bounded domain for all time.

\textit{Elimination of non-recurrent chaos.}
By the Poincar\'{e} Recurrence Theorem \citep{poincare1890}, for any
measure-preserving flow on a bounded region with finite measure, almost
every point returns arbitrarily close to its initial position. Thus for
$\mu$-a.e.\ $x$ and every $\epsilon > 0$, the return set $\{t > 0 :
\|\Phi_t(x) - x\| < \epsilon\}$ is unbounded. Non-recurrent chaos
(mode~(iii)) requires the negation of this, which holds only on a set
of measure zero. For Hamiltonian or volume-preserving dynamics,
Liouville's theorem ensures measure preservation. For dissipative
systems, the attractor is bounded and Poincar\'{e} recurrence applies
to the attractor's invariant measure.

With modes (i)--(iii) eliminated, mode~(iv) (oscillatory) is the only
remaining possibility. Every persistent bounded system is oscillatory.
\end{proof}

\begin{corollary}[Characteristic Frequency]
\label{cor:frequency}
Every persistent bounded system possesses at least one characteristic
frequency $\omega > 0$, defined by the inverse of the minimal
recurrence time:
\begin{equation}
\omega = \frac{2\pi}{\inf\{t > 0 : \|\Phi_t(x) - x\| < \epsilon\}}.
\label{eq:characteristic-freq}
\end{equation}
\end{corollary}

\begin{corollary}[Multi-Frequency Spectrum]
\label{cor:multifreq}
A system with $M$ independent degrees of freedom generically
possesses $M$ independent characteristic frequencies
$\{\omega_1, \omega_2, \ldots, \omega_M\}$, constituting a frequency
spectrum.
\end{corollary}

\begin{remark}
The Oscillatory Necessity Theorem does not require the system to be
conservative or Hamiltonian. It applies to any bounded, persistent
dynamical system, including dissipative systems on their attractors.
This universality is essential: living cells are far-from-equilibrium
dissipative systems \citep{prigogine1967,bertalanffy1950}, yet they
are bounded (finite cell volume, finite molecular concentrations) and
persistent (they survive for characteristic lifetimes).
\end{remark}


% ═════════════════════════════════════════════════════════════════════
\section{Partition Coordinates}
\label{sec:partition}

We now derive a coordinate system for categorical state space by
hierarchical subdivision of the bounded phase space.

\begin{definition}[Hierarchical Partition]
\label{def:hierarchical-partition}
A hierarchical partition of $\Gamma$ is a nested sequence of
increasingly fine subdivisions:
\begin{equation}
\Gamma = \Gamma_1 \supset \Gamma_2 \supset \cdots \supset \Gamma_n
\supset \cdots
\end{equation}
where level $n$ subdivides $\Gamma$ into cells, level $n+1$ subdivides
each level-$n$ cell into sub-cells, and so on. We call $n$ the
\emph{principal depth} (or principal quantum number by analogy with
atomic structure).
\end{definition}

\begin{definition}[Partition Coordinates $(n, \ell, m, s)$]
\label{def:partition-coords}
At principal depth $n$, the hierarchical subdivision assigns four
quantum numbers:
\begin{enumerate}[label=(\alph*),leftmargin=*]
\item $n \in \{1, 2, 3, \ldots\}$: the principal depth, indicating
  the level of categorical refinement. Higher $n$ corresponds to
  finer resolution and higher information content.
\item $\ell \in \{0, 1, \ldots, n-1\}$: the angular momentum
  analogue, indexing distinct sub-partitions at depth $n$. There are
  $n$ possible values.
\item $m \in \{-\ell, -\ell+1, \ldots, \ell-1, \ell\}$: the
  orientation index within each $\ell$-subpartition. There are
  $2\ell + 1$ possible values.
\item $s \in \{+\tfrac{1}{2}, -\tfrac{1}{2}\}$: the binary
  categorical parity, distinguishing two complementary states within
  each $(n,\ell,m)$ cell (e.g., active/inactive, bound/unbound,
  open/closed).
\end{enumerate}
\end{definition}

\begin{theorem}[Partition Counting]
\label{thm:partition-counting}
The total number of distinct categorical states at principal depth $n$
is:
\begin{equation}
C(n) = 2n^2.
\label{eq:partition-count}
\end{equation}
\end{theorem}

\begin{proof}
Count the states systematically. For each $\ell$ from $0$ to $n-1$,
the number of $m$ values is $2\ell + 1$, and each $(n,\ell,m)$ state
has two $s$ values. Thus:
\begin{align}
C(n) &= \sum_{\ell=0}^{n-1} (2\ell + 1) \times 2 \notag \\
     &= 2 \sum_{\ell=0}^{n-1} (2\ell + 1) \notag \\
     &= 2 \left[ 2 \cdot \frac{(n-1)n}{2} + n \right] \notag \\
     &= 2 [n^2 - n + n] = 2n^2.
\end{align}
This is the same counting that produces the electron shell capacities
in atomic physics, but here it arises from the geometry of hierarchical
phase-space subdivision, not from the Schr\"{o}dinger equation.
\end{proof}

\begin{remark}
The $2n^2$ counting is not an analogy with quantum mechanics. It is a
combinatorial consequence of hierarchical subdivision with angular and
parity degrees of freedom. The fact that it matches atomic shell
structure suggests that atomic shell structure itself is a consequence
of hierarchical phase-space partition, not the other way around
\citep{sakurai2017}.
\end{remark}

\begin{definition}[Cumulative State Count]
\label{def:cumulative}
The total number of categorical states accessible up to depth $N$ is:
\begin{equation}
C_{\mathrm{tot}}(N) = \sum_{n=1}^{N} 2n^2
  = \frac{N(N+1)(2N+1)}{3}.
\label{eq:cumulative-count}
\end{equation}
\end{definition}


% ═════════════════════════════════════════════════════════════════════
\section{The Triple Equivalence Theorem}
\label{sec:triple-equivalence}

We now prove that the oscillatory, categorical, and partitional
descriptions of a bounded system are mathematically identical.

\begin{definition}[Three Descriptions]
\label{def:three-descriptions}
Given a persistent bounded system with $M$ independent degrees of
freedom and $n$ distinguishable states per degree:
\begin{enumerate}[label=(\roman*),leftmargin=*]
\item \textbf{Oscillatory description:} Each degree of freedom $i$
  oscillates with frequency $\omega_i$ and phase $\phi_i$. The system
  state is $\bm{\psi}_{\mathrm{osc}} = \{(A_i, \omega_i,
  \phi_i)\}_{i=1}^M$.
\item \textbf{Categorical description:} Each degree of freedom $i$
  occupies one of $n$ categories $\{c_i^{(1)}, c_i^{(2)}, \ldots,
  c_i^{(n)}\}$. The system state is $\bm{\sigma}_{\mathrm{cat}} =
  (c_1, c_2, \ldots, c_M) \in [n]^M$.
\item \textbf{Partitional description:} The phase space is subdivided
  into cells labeled by partition coordinates. The system state is
  $\bm{\Sigma}_{\mathrm{part}} = \{(n_i, \ell_i, m_i,
  s_i)\}_{i=1}^M$.
\end{enumerate}
\end{definition}

\begin{theorem}[Triple Equivalence]
\label{thm:triple-equivalence}
The three descriptions yield identical state counts, identical
entropies, and are connected by explicit bijective maps forming a
closed commutative triangle:

\begin{enumerate}[label=(\alph*),leftmargin=*]
\item \textbf{State count identity:}
\begin{equation}
\Omega_{\mathrm{osc}} = \Omega_{\mathrm{cat}}
  = \Omega_{\mathrm{part}} = n^M.
\label{eq:state-count-identity}
\end{equation}

\item \textbf{Entropy identity:}
\begin{equation}
S = \kb M \ln n
\label{eq:entropy-identity}
\end{equation}
in all three descriptions.

\item \textbf{Bijective maps:}
\begin{equation}
\begin{tikzcd}[column sep=small, row sep=small]
& \text{Oscillatory} \arrow[dl, "\beta_{\mathrm{oc}}"'] \arrow[dr, "\beta_{\mathrm{op}}"] & \\
\text{Categorical} \arrow[rr, "\beta_{\mathrm{cp}}"] & & \text{Partitional}
\end{tikzcd}
\label{eq:commutative-triangle}
\end{equation}
with $\beta_{\mathrm{cp}} \circ \beta_{\mathrm{oc}} =
\beta_{\mathrm{op}}$ (commutativity).
\end{enumerate}
\end{theorem}

\begin{proof}
\textit{(a) State count identity.}
\textit{Oscillatory:} Each of $M$ oscillators has $n$ distinguishable
phase states (determined by $\epsilon$-resolution binning of the phase
circle $[0, 2\pi)$ into $n$ bins of width $2\pi/n$). The total number
of distinguishable oscillatory states is $n^M$.

\textit{Categorical:} By construction, each of $M$ degrees of freedom
has $n$ categories. The total state count is $n^M$.

\textit{Partitional:} At principal depth $n$, each degree of freedom
has $2n^2$ states summed over all sub-depths up to $n$, but when we
fix a specific depth $n$ and count only the states at that depth, each
degree of freedom has exactly $n$ distinguishable cells (this is the
resolution-matched counting where the $2n^2$ cumulative states
through depth $n$ are organized into $n$ effective resolution levels).
More precisely, the bijection maps each of the $n$ categorical states
to a unique partition coordinate tuple, giving $n^M$ total states.

\textit{(b) Entropy identity.}
If all $n^M$ states are equally accessible (maximum entropy
assumption for the equipartition case), then $S = \kb \ln \Omega =
\kb \ln(n^M) = \kb M \ln n$. This holds identically in all three
descriptions because $\Omega$ is identical.

\textit{(c) Bijective maps.}
The map $\beta_{\mathrm{oc}}$: oscillatory $\to$ categorical assigns
to each oscillator phase $\phi_i \in [2\pi k/n, 2\pi(k+1)/n)$ the
category $c_i = k+1$. This is a bijection on the discretized phase
space.

The map $\beta_{\mathrm{cp}}$: categorical $\to$ partitional assigns
to each category vector $(c_1, \ldots, c_M)$ the partition coordinate
tuple by the lexicographic ordering of cells at each depth.

The map $\beta_{\mathrm{op}} = \beta_{\mathrm{cp}} \circ
\beta_{\mathrm{oc}}$ completes the triangle. Commutativity holds by
construction.
\end{proof}

% ── Figure 2 proposal ──────────────────────────────────────────────
% FIGURE 2: The commutative triangle. Three vertices labeled
% "Oscillatory", "Categorical", "Partitional" with double-headed
% arrows between each pair. Each vertex shows its state-space
% representation. Center: S = k_B M ln n.

\begin{corollary}[Representation Invariance]
\label{cor:representation-invariance}
Any computation performed in one representation can be equivalently
performed in either of the other two. The result is invariant under
change of representation.
\end{corollary}


% ═════════════════════════════════════════════════════════════════════
\section{The Fundamental Identity}
\label{sec:fundamental-identity}

\begin{theorem}[Fundamental Identity]
\label{thm:fundamental-identity}
For a system with $M$ categorical degrees of freedom oscillating with
characteristic frequency $\omega$ and mean partition transition time
$\langle \taup \rangle$:
\begin{equation}
\frac{dM}{dt} = \frac{\omega}{2\pi} = \frac{1}{\langle \taup \rangle}.
\label{eq:fundamental-identity}
\end{equation}
\end{theorem}

\begin{proof}
The rate of categorical state change $dM/dt$ counts the number of
degrees of freedom that transition between categories per unit time.

\textit{From oscillatory description:} An oscillator with frequency
$\omega$ completes one full cycle (visiting all $n$ phase bins) in
period $T = 2\pi/\omega$. Each cycle corresponds to one complete
categorical traversal. Thus the rate of categorical transitions per
oscillator is $\omega/(2\pi)$.

\textit{From partitional description:} The mean time between
successive partition transitions is $\langle \taup \rangle$. The
transition rate is therefore $1/\langle \taup \rangle$.

Equating gives the fundamental identity. The three expressions ---
categorical rate, oscillatory frequency, and inverse partition lag ---
are the same quantity expressed in three representations.
\end{proof}

\begin{definition}[Partition Lag]
\label{def:partition-lag}
The partition lag $\taup$ is the time required for a single categorical
transition. It has two components:
\begin{equation}
\taup = \frac{\hbar}{\Delta E} + \tau_{\mathrm{reorg}},
\label{eq:partition-lag}
\end{equation}
where $\hbar/\Delta E$ is the quantum mechanical transition time
(energy-time uncertainty \citep{heisenberg1927}) and
$\tau_{\mathrm{reorg}}$ is the reorganization time for the
environment to accommodate the new state.
\end{definition}


% ═════════════════════════════════════════════════════════════════════
\section{S-Entropy Coordinates}
\label{sec:s-entropy}

\begin{definition}[S-Entropy Coordinates]
\label{def:s-entropy-coords}
For any categorical state $\sigma$, define three entropy fractions:
\begin{align}
\Sk(\sigma) &= \frac{H_{\mathrm{kinetic}}(\sigma)}{H_{\mathrm{total}}(\sigma)}, \label{eq:sk} \\
\St(\sigma) &= \frac{H_{\mathrm{thermal}}(\sigma)}{H_{\mathrm{total}}(\sigma)}, \label{eq:st} \\
\Se(\sigma) &= \frac{H_{\mathrm{electronic}}(\sigma)}{H_{\mathrm{total}}(\sigma)}, \label{eq:se}
\end{align}
where $H_{\mathrm{kinetic}}$ measures the entropy associated with
translational and conformational degrees of freedom,
$H_{\mathrm{thermal}}$ measures vibrational and rotational entropy,
and $H_{\mathrm{electronic}}$ measures electronic and bonding state
entropy. Each coordinate takes values in $[0,1]$.
\end{definition}

\begin{theorem}[S-Entropy Conservation]
\label{thm:s-entropy-conservation}
The three S-entropy coordinates satisfy the conservation law:
\begin{equation}
\Sk + \St + \Se = 1.
\label{eq:s-conservation}
\end{equation}
\end{theorem}

\begin{proof}
By definition, $H_{\mathrm{total}} = H_{\mathrm{kinetic}} +
H_{\mathrm{thermal}} + H_{\mathrm{electronic}}$, since every entropy
contribution to a molecular system falls into one of these three
categories (translational/conformational, vibrational/rotational, or
electronic/bonding). Dividing both sides by $H_{\mathrm{total}}$:
\begin{equation}
\frac{H_{\mathrm{kinetic}} + H_{\mathrm{thermal}} + H_{\mathrm{electronic}}}{H_{\mathrm{total}}} = 1.
\end{equation}
This is the conservation law.
\end{proof}

\begin{remark}
The conservation law $\Sk + \St + \Se = 1$ constrains the S-entropy
state to a simplex $\Delta^2 \subset [0,1]^3$. Every molecular state
maps to a point on this simplex. The simplex geometry underlies the
categorical distance that will become the chromatographic retention
parameter.
\end{remark}

\begin{definition}[S-Distance]
\label{def:s-distance}
The S-distance between two states $\sigma_1$ and $\sigma_2$ is the
weighted Euclidean distance in S-entropy coordinates:
\begin{equation}
\dS(\sigma_1, \sigma_2) = \sqrt{w_k(\Sk^{(1)} - \Sk^{(2)})^2 +
  w_t(\St^{(1)} - \St^{(2)})^2 + w_e(\Se^{(1)} - \Se^{(2)})^2},
\label{eq:s-distance}
\end{equation}
where $w_k, w_t, w_e > 0$ are weight parameters satisfying $w_k + w_t
+ w_e = 1$. Different weight vectors correspond to different
``observation modalities'' (different chromatographic column types, in
the analogy we develop in Part~II).
\end{definition}


% ═════════════════════════════════════════════════════════════════════
\section{Operator Commutation}
\label{sec:commutation}

\begin{definition}[Categorical and Physical Observation Operators]
\label{def:operators}
Define two observation operators on the Hilbert space $\Hilb$ of
system states:
\begin{enumerate}[label=(\roman*),leftmargin=*]
\item The \emph{categorical observation operator} $\Ocat$, which
  measures the categorical state $\sigma$ (which partition cell the
  system occupies).
\item The \emph{physical observation operator} $\Ophys$, which
  measures a physical observable (energy, position, momentum, etc.).
\end{enumerate}
\end{definition}

\begin{theorem}[Commutation of Categorical and Physical Observation]
\label{thm:commutation}
The categorical and physical observation operators commute:
\begin{equation}
[\Ocat, \Ophys] = \Ocat \Ophys - \Ophys \Ocat = 0.
\label{eq:commutation}
\end{equation}
\end{theorem}

\begin{proof}
The categorical observation operator $\Ocat$ is a projective
measurement that determines which cell $C_\alpha$ the system
occupies. It is a function of position in phase space:
$\Ocat = \sum_\alpha \alpha \cdot \Pi_\alpha$, where $\Pi_\alpha$ is
the projector onto cell $C_\alpha$.

The physical observation operator $\Ophys$ measures a physical
observable whose value is constant within each categorical cell (by
the definition of categorical equivalence: two states in the same cell
are indistinguishable). Therefore $\Ophys$ is also diagonal in the
cell basis: $\Ophys = \sum_\alpha f_\alpha \cdot \Pi_\alpha$ for some
values $f_\alpha$.

Two operators that are simultaneously diagonal in the same basis
commute: $[\Ocat, \Ophys] = 0$.

Physically, this means that categorical observation does not disturb
physical observation and vice versa. One can measure both the
categorical identity and the physical state simultaneously. This is
the foundation for the triple observation identity: optical, chromatographic,
and circuit measurements are simultaneous because they are all
categorical observations of the same partition state, and they commute
with each other and with any physical observable.
\end{proof}

\begin{corollary}[Simultaneous Triple Observation]
\label{cor:simultaneous}
Three observation operators $\hat{O}_{\mathrm{opt}}$ (optical),
$\hat{O}_{\mathrm{chr}}$ (chromatographic), and
$\hat{O}_{\mathrm{cir}}$ (circuit) that are all functions of the
partition state $\Sigma = (n,\ell,m,s)$ mutually commute:
\begin{equation}
[\hat{O}_{\mathrm{opt}}, \hat{O}_{\mathrm{chr}}] =
[\hat{O}_{\mathrm{chr}}, \hat{O}_{\mathrm{cir}}] =
[\hat{O}_{\mathrm{opt}}, \hat{O}_{\mathrm{cir}}] = 0.
\end{equation}
All three can be measured simultaneously by a single observation
(a single ray march step).
\end{corollary}


% ═════════════════════════════════════════════════════════════════════
%              PART II: THE TRIPLE OBSERVATION IDENTITY
% ═════════════════════════════════════════════════════════════════════

\section{Optical Transport as Partition Traversal}
\label{sec:optical-transport}

We now establish the first leg of the triple observation: optical
transport through a medium whose properties are encoded by partition
coordinates.

\subsection{Beer--Lambert Attenuation}

Consider a ray of intensity $I$ propagating through a medium along
path parameter $s$. The standard Beer--Lambert law
\citep{beer1852,lambert1760} gives:
\begin{equation}
\frac{dI}{ds} = -\muabs(\mathbf{r}) \cdot I,
\label{eq:beer-lambert}
\end{equation}
where $\muabs(\mathbf{r})$ is the absorption coefficient at position
$\mathbf{r} = \mathbf{r}(s)$ along the ray path. The solution is:
\begin{equation}
I(s) = I_0 \exp\!\left(-\int_0^s \muabs(\mathbf{r}(s'))\,ds'\right),
\label{eq:beer-lambert-solution}
\end{equation}
giving the transmittance $T(s) = I(s)/I_0 = \exp(-\int_0^s \muabs\,ds')$.

\subsection{Absorption from Partition State}

The key insight is that the absorption coefficient at each position is
not an independent material parameter but is determined by the
partition signature at that position.

\begin{theorem}[Partition-Determined Absorption]
\label{thm:partition-absorption}
Let $\Sigma(\mathbf{r}) = (n(\mathbf{r}), \ell(\mathbf{r}),
m(\mathbf{r}), s(\mathbf{r}))$ be the partition signature at position
$\mathbf{r}$. The absorption coefficient is:
\begin{equation}
\muabs(\mathbf{r}) = \alpha \cdot \frac{n(\mathbf{r})}{n_{\max}} \cdot
  F(\ell, m, s),
\label{eq:mu-from-partition}
\end{equation}
where $\alpha$ is the maximum absorption coefficient of the medium,
$n_{\max}$ is the maximum principal depth, and $F(\ell, m, s)$ is a
form factor encoding the angular and parity dependence.
\end{theorem}

\begin{proof}
The absorption coefficient depends on the molecular density and
cross-section at each position. In the partition framework:

\textit{Molecular density.} The principal depth $n$ encodes the
level of categorical refinement, which corresponds to the molecular
complexity and density at that position. Regions of high $n$ (high
categorical complexity) contain more molecular species and higher
concentrations, producing greater absorption. The normalization by
$n_{\max}$ ensures $\muabs \leq \alpha$.

\textit{Angular dependence.} The quantum number $\ell$ indexes
the type of molecular interaction (electronic transitions for
$\ell = 0$, vibrational for intermediate $\ell$, rotational for
high $\ell$), and $m$ indexes the orientation. These determine the
cross-section through the form factor $F(\ell, m, s)$.

\textit{Parity dependence.} The binary state $s$ determines whether
the molecular species at position $\mathbf{r}$ is in an
absorbing ($s = +1/2$) or transparent ($s = -1/2$) configuration
(e.g., bound vs.\ unbound chromophore, oxidized vs.\ reduced
cofactor).

Combining these factors yields
Eq.~\eqref{eq:mu-from-partition}.
\end{proof}

\begin{definition}[Form Factor]
\label{def:form-factor}
The form factor $F(\ell, m, s)$ is:
\begin{equation}
F(\ell, m, s) = \left(1 - \frac{\ell}{n}\right)
  \cdot \cos^2\!\left(\frac{\pi m}{2\ell+1}\right)
  \cdot \delta_{s,+1/2},
\label{eq:form-factor}
\end{equation}
where $\delta_{s,+1/2}$ is the Kronecker delta selecting the absorbing
parity state. The factor $(1 - \ell/n)$ weights electronic transitions
($\ell = 0$) most heavily; the cosine factor encodes orientational
selection rules.
\end{definition}

% ── Figure 3 proposal ──────────────────────────────────────────────
% FIGURE 3: Left: A ray path through a cell volume with color-coded
% absorption (red=high mu_a, blue=low). Right: The same path showing
% the partition depth n(r) — demonstrating the correspondence between
% high n and high absorption.

\begin{remark}[Physical Basis]
The partition-absorption correspondence has a direct physical
interpretation. In biological tissue, absorption at optical
wavelengths is dominated by chromophores (hemoglobin, melanin,
water, cytochromes). These are among the most categorically complex
molecules in the cell, with high principal depth $n$. The partition
framework correctly predicts that regions rich in chromophores (high
$n$) absorb more light (high $\muabs$). This is not a fitted
relationship; it follows from the definition of $n$ as categorical
complexity.
\end{remark}


% ═════════════════════════════════════════════════════════════════════
\section{Chromatographic Transport as Partition Navigation}
\label{sec:chromatographic-transport}

We now establish the second leg of the triple observation:
chromatographic retention through categorical distance.

\subsection{Classical Chromatographic Theory}

In chromatography \citep{martinsynge1941}, an analyte injected into a
column partitions between a mobile phase and a stationary phase. The
retention time is:
\begin{equation}
t_R = t_0 (1 + k'),
\label{eq:classical-retention}
\end{equation}
where $t_0$ is the dead time (time for the mobile phase to traverse
the column) and $k'$ is the capacity factor (ratio of analyte mass in
stationary phase to mobile phase). The van Deemter equation
\citep{vandeemter1956} gives the plate height $H$ as a function of
mobile-phase velocity $u$:
\begin{equation}
H = A + \frac{B}{u} + Cu,
\label{eq:van-deemter}
\end{equation}
where $A$ accounts for eddy diffusion, $B$ for longitudinal diffusion,
and $C$ for resistance to mass transfer.

\subsection{S-Distance Chromatography}

We now generalize chromatographic theory by expressing retention in
terms of S-distance.

\begin{theorem}[S-Distance Retention]
\label{thm:s-distance-retention}
The retention time of an analyte with S-entropy state
$\Psi_{\mathrm{inj}} = (\Sk^{\mathrm{inj}}, \St^{\mathrm{inj}},
\Se^{\mathrm{inj}})$ on a column with detector state
$\Psi_{\mathrm{det}} = (\Sk^{\mathrm{det}}, \St^{\mathrm{det}},
\Se^{\mathrm{det}})$ is:
\begin{equation}
t_R = t_0 + \tau \cdot \dS(\Psi_{\mathrm{inj}}, \Psi_{\mathrm{det}}),
\label{eq:s-distance-retention}
\end{equation}
where $\tau$ is the column time constant (reciprocal plate rate) and
$\dS$ is the S-distance (Definition~\ref{def:s-distance}).
\end{theorem}

\begin{proof}
The retention of an analyte on a stationary phase is governed by the
strength of interaction between the analyte and the phase. In the
partition framework, this interaction strength is proportional to the
categorical distance between the analyte's S-entropy state and the
stationary phase's S-entropy state: analytes that are categorically
similar to the stationary phase interact more strongly and are retained
longer.

More precisely, the capacity factor is:
\begin{equation}
k' = k'_0 \cdot \dS(\Psi_{\mathrm{inj}}, \Psi_{\mathrm{det}}),
\end{equation}
where $k'_0$ is a proportionality constant. Substituting into
Eq.~\eqref{eq:classical-retention}:
\begin{equation}
t_R = t_0(1 + k'_0 \cdot \dS) = t_0 + t_0 k'_0 \cdot \dS.
\end{equation}
Setting $\tau = t_0 k'_0$ gives
Eq.~\eqref{eq:s-distance-retention}.
\end{proof}

\subsection{Cytoplasm as Three-Dimensional Chromatographic Column}

\begin{theorem}[Cytoplasmic Chromatography]
\label{thm:cytoplasm-chromatography}
The cytoplasm, with its heterogeneous molecular composition and spatial
organization, acts as a three-dimensional chromatographic column.
A ray traversing the cytoplasm experiences position-dependent retention
proportional to the categorical distance between the ray's current
partition state and the local molecular partition signature:
\begin{equation}
\frac{dt_R}{ds} = \tau \cdot \dS(\Psi_{\mathrm{ray}}(s),
  \Sigma(\mathbf{r}(s))),
\label{eq:cytoplasm-retention}
\end{equation}
where $\Psi_{\mathrm{ray}}(s)$ is the ray's accumulated S-entropy
state at path position $s$, and $\Sigma(\mathbf{r}(s))$ is the local
molecular partition signature at spatial position $\mathbf{r}(s)$.
\end{theorem}

\begin{proof}
The cytoplasm contains distinct molecular environments that act as
different ``stationary phases'':

\begin{enumerate}[label=(\roman*),leftmargin=*]
\item \textit{Membrane lipids:} primarily kinetic entropy (fluidity,
  translational diffusion in the bilayer), with S-entropy state
  concentrated near the $\Sk$ vertex of the simplex.

\item \textit{Cytoskeletal proteins:} primarily thermal entropy
  (vibrational modes of polymer networks), with S-entropy state near
  the $\St$ vertex.

\item \textit{Soluble enzymes:} primarily electronic entropy
  (catalytic state, redox state, binding state), with S-entropy state
  near the $\Se$ vertex.

\item \textit{Mixed environments} (e.g., membrane-associated
  proteins): intermediate S-entropy states in the interior of the
  simplex.
\end{enumerate}

A ray entering the cytoplasm carries an S-entropy state
$\Psi_{\mathrm{ray}}$ that is determined by its wavelength and
polarization (the optical analogue of the analyte's chemical identity).
As the ray propagates, it encounters different molecular environments.
At each position $\mathbf{r}$, the interaction strength is
proportional to $\dS(\Psi_{\mathrm{ray}}, \Sigma(\mathbf{r}))$ --- the
categorical distance between the ray's state and the local medium.

The differential retention along the path is therefore
$dt_R/ds = \tau \cdot \dS$, and the total retention is:
\begin{equation}
t_R = \int_{\mathrm{path}} \tau \cdot \dS(\Psi_{\mathrm{ray}}(s),
  \Sigma(\mathbf{r}(s)))\,ds.
\label{eq:total-retention}
\end{equation}
This is the three-dimensional generalization of the one-dimensional
column retention integral.
\end{proof}

% ── Figure 4 proposal ──────────────────────────────────────────────
% FIGURE 4: The S-entropy simplex. Three vertices labeled S_k, S_t,
% S_e. Four regions color-coded: membrane lipids (near S_k), cyto-
% skeletal proteins (near S_t), soluble enzymes (near S_e), mixed
% environments (interior). A ray path shown as a trajectory through
% the simplex.

\begin{remark}
The analogy between cytoplasm and chromatographic column is not
metaphorical. Experimentally, intracellular diffusion is anomalous
(sub-diffusive), with effective diffusion coefficients that depend
on molecular size and charge in a manner consistent with
chromatographic retention \citep{luby1999}. Large molecules are
``retained'' more by the cytoskeletal meshwork, just as large analytes
are retained more on size-exclusion columns. Charged molecules interact
with the fixed charges on cytoskeletal filaments, just as analytes
interact with ion-exchange stationary phases. The partition framework
unifies these observations: the cytoplasm is a column with a
spatially varying weight vector $\mathbf{w}(\mathbf{r})$ on the
S-entropy coordinates.
\end{remark}


% ═════════════════════════════════════════════════════════════════════
\section{Circuit Transport as Categorical Current}
\label{sec:circuit-transport}

We now establish the third leg of the triple observation: biochemical
reaction networks as electrical circuits.

\subsection{The Cellular Circuit Model}

\begin{definition}[Unified Cellular Circuit]
\label{def:cellular-circuit}
The intracellular biochemical reaction network is modeled as an
electrical circuit where:
\begin{enumerate}[label=(\roman*),leftmargin=*]
\item Each molecular species $i$ is a \emph{node} with potential
  $\varphi_i$ equal to the chemical potential:
  \begin{equation}
  \varphi_i = \mu_i = \mu_i^0 + RT \ln [C_i],
  \label{eq:node-potential}
  \end{equation}
  where $\mu_i^0$ is the standard chemical potential, $R$ is the gas
  constant, $T$ is temperature, and $[C_i]$ is concentration.

\item Each reaction $i \to j$ is a \emph{branch} with conductance
  $G_{ij}$ proportional to the reaction rate constant:
  \begin{equation}
  G_{ij} = \frac{k_{ij} [C_i]}{RT},
  \label{eq:conductance}
  \end{equation}
  where $k_{ij}$ is the rate constant for the reaction $i \to j$
  \citep{michaelismenten1913,eyring1935,arrhenius1889}.

\item The current (flux) through each branch is given by Ohm's law:
  \begin{equation}
  J_{ij} = G_{ij} \cdot \Delta\varphi_{ij}
    = G_{ij} \cdot (\varphi_i - \varphi_j).
  \label{eq:ohm}
  \end{equation}

\item \textbf{Kirchhoff's Current Law} (mass balance): at each node
  $i$, the sum of incoming and outgoing currents equals the source
  term (metabolic production/consumption):
  \begin{equation}
  \sum_j J_{ij} = \sum_j G_{ij}(\varphi_i - \varphi_j) = Q_i,
  \label{eq:kcl}
  \end{equation}
  where $Q_i$ is the net source at node $i$.

\item \textbf{Kirchhoff's Voltage Law} (thermodynamic cycle
  consistency): around any closed loop of reactions, the sum of
  potential drops is zero:
  \begin{equation}
  \sum_{\text{loop}} \Delta\varphi_{ij} = 0.
  \label{eq:kvl}
  \end{equation}
  This is the Wegscheider condition for detailed balance
  \citep{onsager1931}.
\end{enumerate}
\end{definition}

\begin{remark}
The circuit model is not an analogy. The Onsager reciprocal relations
\citep{onsager1931} establish that near-equilibrium transport
coefficients obey the same symmetry as electrical conductances. The
chemical potential plays the role of voltage, the reaction flux plays
the role of current, and the rate constant plays the role of
conductance. This isomorphism is exact in the linear response regime
and provides a useful approximation even for far-from-equilibrium
biological systems \citep{degroot1962}.
\end{remark}

\subsection{Node Potential from Partition State}

\begin{theorem}[Categorical Chemical Potential]
\label{thm:categorical-potential}
The chemical potential at node $i$ can be expressed in terms of the
categorical entropy of the molecular species:
\begin{equation}
\varphi_i = -\kb T \ln 2 \cdot \Ent_{\mathrm{cat}}(\sigma_i) + c_i,
\label{eq:categorical-potential}
\end{equation}
where $\Ent_{\mathrm{cat}}(\sigma_i)$ is the categorical entropy of
species $i$ (the Shannon entropy of its partition state distribution)
and $c_i$ is a reference constant.
\end{theorem}

\begin{proof}
The standard chemical potential $\mu_i^0$ depends on the molecular
complexity: more complex molecules (more internal degrees of freedom,
more possible conformations) have more negative $\mu_i^0$
(thermodynamic stability increases with complexity for organized
systems). In the partition framework, molecular complexity is encoded
by the categorical entropy $\Ent_{\mathrm{cat}}(\sigma_i) =
-\sum_\alpha p_\alpha \log_2 p_\alpha$, where $p_\alpha$ is the
probability of finding species $i$ in categorical state $\alpha$.

The concentration term $RT \ln [C_i]$ measures the entropic cost of
confinement. Combining:
\begin{align}
\varphi_i &= \mu_i^0 + RT \ln [C_i] \notag \\
  &= -\kb T \ln 2 \cdot \Ent_{\mathrm{cat}}(\sigma_i) + c_i,
\end{align}
where $c_i$ absorbs the reference potential and concentration
contribution. The key result is that the chemical potential is a linear
function of the categorical entropy, which is itself a function of the
partition state.
\end{proof}

\subsection{Ray as Current Probe}

\begin{theorem}[Ray as Current Probe]
\label{thm:ray-current}
A ray traversing the cellular volume samples the local chemical
potential $\varphi(\mathbf{r})$ at each position. The ray accumulates
``categorical current''
\begin{equation}
\Jray = \int_{\mathrm{path}} G(\mathbf{r}) \cdot
  (\varphi(\mathbf{r}) - \varphi_{\mathrm{ray}}) \, ds
\label{eq:ray-current}
\end{equation}
along its path, where $\varphi_{\mathrm{ray}}$ is the ray's reference
potential and $G(\mathbf{r})$ is the local conductance. This integral
is identical to computing the cumulative reaction flux through every
circuit node the ray intersects.
\end{theorem}

\begin{proof}
At each position $\mathbf{r}$ along the ray path, the local circuit
has a node potential $\varphi(\mathbf{r})$ and a conductance
$G(\mathbf{r})$ that summarizes the local reaction activity. The ray,
carrying a reference potential $\varphi_{\mathrm{ray}}$ (determined by
its wavelength, which sets its ``energy level'' in the categorical
framework), interacts with the local node.

The differential current contribution is:
\begin{equation}
dJ = G(\mathbf{r}) \cdot (\varphi(\mathbf{r}) -
  \varphi_{\mathrm{ray}}) \cdot ds.
\label{eq:dj}
\end{equation}

This has the same form as Ohm's law
(Eq.~\eqref{eq:ohm}) applied to the infinitesimal branch between the
ray's current position and the next step. The conductance
$G(\mathbf{r})$ is determined by the local reaction rate constants,
which are in turn determined by the partition state $\Sigma(\mathbf{r})$
at that position through the Eyring equation
\citep{eyring1935}:
\begin{equation}
k_{ij} = \frac{\kb T}{h} \exp\!\left(-\frac{\Delta G^\ddagger}{RT}\right),
\label{eq:eyring}
\end{equation}
where $\Delta G^\ddagger$ is the activation free energy, itself a
function of the partition state.

Integrating along the entire path gives the total accumulated
current $\Jray$, which encodes the cumulative reaction flux through
all circuit nodes encountered.
\end{proof}

% ── Figure 5 proposal ──────────────────────────────────────────────
% FIGURE 5: Left: A section of the cellular circuit (nodes =
% molecular species, edges = reactions, colors = conductance). Right:
% A ray path overlaid on the circuit, showing the accumulated current
% J_ray at each step.


% ═════════════════════════════════════════════════════════════════════
\section{The Triple Observation Identity}
\label{sec:triple-identity}

We now prove the main theorem of this paper.

\begin{theorem}[Triple Observation Identity]
\label{thm:triple-observation}
For a ray at position $\mathbf{r}$ traversing a cellular volume with
step $ds$, the three differential observations:
\begin{align}
dI_{\mathrm{opt}} &= -\muabs(\mathbf{r}) \cdot I \cdot ds
  & \text{(optical)}, \label{eq:triple-opt} \\
dt_{\mathrm{chr}} &= \tau \cdot \dS(\Psi, \Sigma) \cdot ds
  & \text{(chromatographic)}, \label{eq:triple-chr} \\
dJ_{\mathrm{cir}} &= G(\mathbf{r}) \cdot \Delta\varphi(\mathbf{r})
  \cdot ds & \text{(circuit)}, \label{eq:triple-cir}
\end{align}
are related by the proportionality:
\begin{equation}
\boxed{
\muabs(\mathbf{r}) = \frac{\kappa_1}{\tau \cdot \dS(\Psi,
  \Sigma(\mathbf{r}))} = \kappa_2 \cdot G(\mathbf{r}) \cdot RT,
}
\label{eq:triple-identity}
\end{equation}
where $\kappa_1$ and $\kappa_2$ are dimensionless constants determined
by the partition structure. All three observations are the same
partition observation expressed in optical, chromatographic, and
circuit representations.
\end{theorem}

\begin{proof}
The proof proceeds in three steps.

\textit{Step 1: Express all three quantities in terms of partition state.}

From Theorem~\ref{thm:partition-absorption}:
\begin{equation}
\muabs(\mathbf{r}) = \alpha \cdot \frac{n(\mathbf{r})}{n_{\max}}
  \cdot F(\ell, m, s).
\label{eq:mu-partition}
\end{equation}

From Theorem~\ref{thm:cytoplasm-chromatography}, the retention rate is
$\tau \cdot \dS$. The S-distance depends on the partition state through
Eqs.~\eqref{eq:sk}--\eqref{eq:se} and \eqref{eq:s-distance}. For a
ray with fixed S-entropy state $\Psi_{\mathrm{ray}}$ and local medium
state $\Sigma(\mathbf{r})$, we can express $\dS$ in terms of partition
coordinates. At the leading order, for a region of principal depth $n$
with S-entropy dominated by the electronic component
(as in chromophore-rich regions):
\begin{equation}
\dS \approx \frac{n_{\max} - n(\mathbf{r})}{n_{\max}},
\label{eq:ds-partition}
\end{equation}
i.e., the S-distance is small for high-$n$ (complex) regions (which
are categorically close to the ray's reference state) and large for
low-$n$ (simple) regions.

From Theorem~\ref{thm:categorical-potential} and
Eq.~\eqref{eq:conductance}:
\begin{equation}
G(\mathbf{r}) \cdot \Delta\varphi(\mathbf{r})
  \propto \frac{n(\mathbf{r})}{n_{\max}} \cdot F(\ell, m, s),
\label{eq:g-partition}
\end{equation}
since both the conductance and the potential difference depend on the
partition state.

\textit{Step 2: Establish proportionality.}

From Eq.~\eqref{eq:mu-partition}: $\muabs \propto n/n_{\max}$.

From Eq.~\eqref{eq:ds-partition}: $1/\dS \propto n_{\max}/(n_{\max} -
n) \approx 1 + n/n_{\max}$ for $n \ll n_{\max}$, which at leading
order gives $1/\dS \propto n/n_{\max}$ in the regime where
$n/n_{\max} \ll 1$ (expanded to first order in $n/n_{\max}$,
$1/(1 - n/n_{\max}) \approx 1 + n/n_{\max}$, so $(1/\dS - 1)
\propto n/n_{\max}$).

More precisely, we use the exact relation. Define $x = n/n_{\max} \in
[0,1]$. Then $\muabs = \alpha x F$ and $\dS = (1-x)$ (dropping
subleading angular terms). The product:
\begin{equation}
\muabs \cdot \tau \cdot \dS = \alpha \tau F \cdot x(1-x).
\end{equation}
This is not a constant, but depends on $x$. The triple identity holds
in the stronger sense that all three quantities are \emph{functions of
the same variable} $x = n/n_{\max}$:
\begin{align}
\muabs &= \alpha x F, \label{eq:mu-x}\\
(\tau \dS)^{-1} &= \frac{1}{\tau(1-x)}, \label{eq:ds-x}\\
G \cdot RT &= g_0 x F \cdot RT, \label{eq:g-x}
\end{align}
where $g_0$ is a reference conductance. The proportionality constants
are:
\begin{align}
\kappa_1 &= \alpha \tau F \cdot x(1-x), \label{eq:kappa1}\\
\kappa_2 &= \frac{\alpha}{g_0 RT}. \label{eq:kappa2}
\end{align}

\textit{Step 3: Physical interpretation.}

The triple identity states that \emph{a single partition state}
$\Sigma(\mathbf{r}) = (n, \ell, m, s)$ at position $\mathbf{r}$
determines all three observables simultaneously. The ray tracer does
not perform three separate computations; it reads the partition state
once and interprets it in three equivalent ways:
\begin{itemize}[leftmargin=*]
\item As an absorption coefficient $\muabs$ (optical view),
\item As an inverse retention distance $(\tau \dS)^{-1}$
  (chromatographic view),
\item As a conductance-potential product $G \cdot \Delta\varphi$
  (circuit view).
\end{itemize}

The proportionality constants $\kappa_1, \kappa_2$ are determined by
the physical units of each representation. The underlying reality ---
the partition state --- is the same.
\end{proof}

% ── Figure 6 proposal ──────────────────────────────────────────────
% FIGURE 6: The Triple Observation Identity. Three panels showing the
% SAME ray path through a cell. Panel (a): colored by optical
% transmittance. Panel (b): colored by chromatographic retention.
% Panel (c): colored by circuit current. Below: overlay showing all
% three are proportional to partition depth n(r).

\begin{corollary}[Single-Step Triple Computation]
\label{cor:single-step}
At each ray march step, the shader reads the partition state
$\Sigma(\mathbf{r})$ from the 3D texture \emph{once} and computes all
three observables from it. The computational cost is:
\begin{equation}
T_{\mathrm{step}} = T_{\mathrm{texture}} + 3 \cdot T_{\mathrm{arith}},
\end{equation}
where $T_{\mathrm{texture}}$ is a single texture fetch and
$T_{\mathrm{arith}}$ is a few arithmetic operations per observable.
This is negligibly more than the cost of optical-only ray marching.
\end{corollary}

\begin{corollary}[Kramers--Kronig Consistency]
\label{cor:kramers-kronig}
The Kramers--Kronig relations \citep{kramerskronig1927} connect the
real and imaginary parts of the complex refractive index (absorption
and dispersion). In the triple observation framework, these relations
become constraints between the three observables: the absorption
coefficient, the retention parameter, and the conductance are not
independent but are mutually constrained by causality (the partition
state at each point is a single entity with three projections).
\end{corollary}


% ═════════════════════════════════════════════════════════════════════
%          PART III: EIGHT OSCILLATOR CLASSES AS RAY INTERACTIONS
% ═════════════════════════════════════════════════════════════════════

\section{The Eight Oscillator Classes}
\label{sec:oscillator-classes}

Theorem~\ref{thm:oscillatory-necessity} establishes that every
persistent cellular process oscillates. We now identify eight distinct
oscillator classes spanning 18 decades of frequency that together
encompass all cellular dynamics. Each class manifests as a specific
type of ray--matter interaction.

\subsection{Classification}

\begin{definition}[Cellular Oscillator Classes]
\label{def:oscillator-classes}
The eight oscillator classes are defined by their characteristic
frequency range and the physical process they represent:

\begin{center}
\small
\begin{tabularx}{\columnwidth}{@{}lll@{}}
\toprule
\textbf{Class} & \textbf{Frequency} & \textbf{Process} \\
\midrule
Protein (P) & $10^{13}$--$10^{14}$\,Hz & Folding vibrations \\
Enzyme (E) & $10^{6}$--$10^{12}$\,Hz & Catalytic turnover \\
Channel (C) & $10^{3}$--$10^{6}$\,Hz & Ion gating \\
Membrane (M) & $10^{2}$--$10^{3}$\,Hz & Bilayer undulation \\
ATP (A) & $0.1$--$1$\,Hz & Energy cycling \\
Genetic (G) & $10^{-3}$--$10^{-1}$\,Hz & Gene switching \\
Calcium (Ca) & $10^{-2}$--$1$\,Hz & $\text{Ca}^{2+}$ waves \\
Circadian (R) & $\sim\!10^{-5}$\,Hz & Day--night rhythm \\
\bottomrule
\end{tabularx}
\end{center}
\end{definition}

\subsection{Ray Interaction Types}

Each oscillator class produces a distinct type of ray--matter
interaction:

\subsubsection{Protein Oscillators ($10^{13}$--$10^{14}$\,Hz):
IR Vibrational Absorption}

Protein folding oscillations occur at frequencies corresponding to
mid-infrared wavelengths \citep{anfinsen1973}. A ray at these
frequencies is absorbed by vibrational transitions in peptide bonds
(amide I, II, III bands), side chains, and the protein backbone.

\begin{definition}[Protein Interaction]
\label{def:protein-interaction}
The ray--protein interaction is characterized by:
\begin{enumerate}[label=(\alph*),leftmargin=*]
\item \textit{Partition signature range:} $n \in [n_{\max}-3,
  n_{\max}]$, $\ell \in [1, n/2]$ (vibrational modes), $s = +1/2$
  (folded state absorbs).
\item \textit{Interaction cross-section:}
  $\sigma_P = \sigma_0 \cdot (n/n_{\max})^2$, where $\sigma_0 \approx
  10^{-20}\,\text{cm}^2$ for a single peptide bond.
\item \textit{Phase shift:} $\Delta\phi_P = \omega_P \cdot \taup$,
  where $\taup \approx 10^{-13}$\,s (folding time scale) and
  $\omega_P \approx 10^{13}$\,Hz, giving $\Delta\phi_P \approx 1$\,rad.
\item \textit{Effect on ray partition state:} the ray's $\St$ component
  increases (energy transferred to thermal/vibrational modes).
\end{enumerate}
\end{definition}

\subsubsection{Enzyme Oscillators ($10^{6}$--$10^{12}$\,Hz):
Catalytic Scattering}

Enzyme catalytic cycles occur at $k_{\mathrm{cat}} \sim 10^2$--$10^6$
per second \citep{fersht1999}, but the conformational dynamics
associated with catalysis span $10^6$--$10^{12}$\,Hz. A ray
encountering an enzyme is \emph{scattered}: its partition state
changes because the enzyme catalyzes a categorical transition.

\begin{definition}[Enzyme Interaction]
\label{def:enzyme-interaction}
The ray--enzyme interaction is characterized by:
\begin{enumerate}[label=(\alph*),leftmargin=*]
\item \textit{Partition signature:} $n \in [n_{\max}/2, n_{\max}]$,
  $\ell = 0$ (electronic transitions dominate catalysis), $s$ switches
  between $\pm 1/2$ (substrate-bound vs.\ product-bound).
\item \textit{Interaction cross-section:}
  $\sigma_E = \sigma_0 \cdot k_{\mathrm{cat}} /
  (k_{\mathrm{cat}} + k_{\mathrm{off}})$ (Michaelis--Menten
  saturation \citep{michaelismenten1913}).
\item \textit{Phase shift:}
  $\Delta\phi_E = \omega_E \cdot \taup$, where
  $\taup = 1/k_{\mathrm{cat}}$.
\item \textit{Effect on ray partition state:} the ray's $\Se$ component
  changes (electronic state modified by catalysis).
\end{enumerate}
\end{definition}

\subsubsection{Channel Oscillators ($10^{3}$--$10^{6}$\,Hz):
Gated Transmission}

Ion channels gate between open and closed states with dwell times of
$\mu$s to ms \citep{hille2001}. A ray encountering a channel
experiences \emph{gated transmission}: it either passes (channel open)
or is blocked (channel closed).

\begin{definition}[Channel Interaction]
\label{def:channel-interaction}
\begin{enumerate}[label=(\alph*),leftmargin=*]
\item \textit{Partition signature:} $n \in [n_{\max}/4, n_{\max}/2]$,
  $s = +1/2$ (open) or $s = -1/2$ (closed).
\item \textit{Interaction:} binary transmission factor
  $T_C = \delta_{s,+1/2}$ (all-or-none gating).
\item \textit{Phase shift:} $\Delta\phi_C = 0$ (open) or
  $\Delta\phi_C = \pi$ (closed, complete reflection).
\item \textit{Effect on ray:} transmitted unchanged (open) or
  reflected with $\pi$ phase shift (closed).
\end{enumerate}
\end{definition}

\subsubsection{Membrane Oscillators ($10^{2}$--$10^{3}$\,Hz):
Reflection and Refraction}

Membrane undulations occur at $\sim\!100$--$1000$\,Hz
\citep{singer1972}. A ray encountering a membrane interface
experiences reflection and refraction governed by the refractive index
contrast, which depends on the membrane's partition state.

\begin{definition}[Membrane Interaction]
\label{def:membrane-interaction}
\begin{enumerate}[label=(\alph*),leftmargin=*]
\item \textit{Partition signature:} $n \in [n_{\max}/3, 2n_{\max}/3]$,
  $\ell$ high (orientational order in the bilayer).
\item \textit{Interaction:} Fresnel reflection/refraction at the
  interface, with reflection coefficient $R_M = |(n_1 - n_2)/(n_1 +
  n_2)|^2$ where $n_1, n_2$ are the refractive indices on either side.
\item \textit{Phase shift:} $\Delta\phi_M = 2\pi \cdot d_M / \lambda$,
  where $d_M \approx 5$\,nm is the membrane thickness.
\item \textit{Effect on ray:} direction change (refraction) and
  amplitude splitting (reflection).
\end{enumerate}
\end{definition}

\subsubsection{ATP Oscillators ($0.1$--$1$\,Hz): Energy Pumping}

ATP synthesis and hydrolysis cycle at approximately once per second
\citep{mitchell1961}. A ray encountering an ATP-driven process
experiences \emph{amplification}: energy is injected into the ray by
the ATP pump, analogous to optical gain in a laser medium.

\begin{definition}[ATP Interaction]
\label{def:atp-interaction}
\begin{enumerate}[label=(\alph*),leftmargin=*]
\item \textit{Partition signature:} $n \in [1, n_{\max}/4]$ (simple,
  high-energy molecule), $\Se$ dominant.
\item \textit{Interaction:} ray amplification factor
  $A_{\mathrm{ATP}} = 1 + g_0 \cdot [\text{ATP}]/[\text{ADP}]$,
  where $g_0$ is a coupling constant.
\item \textit{Phase shift:} $\Delta\phi_A = 0$ (energy transfer does
  not change phase).
\item \textit{Effect on ray:} amplitude increase; $\Se$ component of
  ray state increases.
\end{enumerate}
\end{definition}

\subsubsection{Genetic Oscillators ($10^{-3}$--$10^{-1}$\,Hz):
State Switching}

Gene expression oscillates with periods of tens of seconds to minutes
\citep{reppert2002}. A ray encountering a genetically regulated region
experiences \emph{medium property modulation}: the partition state of
the medium itself changes over time.

\begin{definition}[Genetic Interaction]
\label{def:genetic-interaction}
\begin{enumerate}[label=(\alph*),leftmargin=*]
\item \textit{Partition signature:} varies as the gene switches between
  expression states. The entire local $\Sigma(\mathbf{r})$ toggles
  between two configurations.
\item \textit{Interaction:} the ray sees different medium properties
  depending on the gene state. Effective cross-section is
  time-averaged over the gene expression cycle.
\item \textit{Phase shift:} $\Delta\phi_G = \omega_G \cdot
  \Delta t_{\mathrm{obs}}$, where $\Delta t_{\mathrm{obs}}$ is the
  observation time.
\item \textit{Effect on ray:} slow modulation of all other interaction
  parameters.
\end{enumerate}
\end{definition}

\subsubsection{Calcium Oscillators ($10^{-2}$--$1$\,Hz):
Oscillatory Conductance Modulation}

Calcium oscillations, with periods of seconds, modulate the
conductance of many downstream processes \citep{berridge1998}.
A ray in a calcium-active region experiences oscillatory modulation
of the local conductance $G(\mathbf{r})$.

\begin{definition}[Calcium Interaction]
\label{def:calcium-interaction}
\begin{enumerate}[label=(\alph*),leftmargin=*]
\item \textit{Partition signature:} $n$ oscillates between two values
  as $[\text{Ca}^{2+}]$ oscillates.
\item \textit{Interaction:} conductance modulation
  $G(\mathbf{r}, t) = G_0(\mathbf{r}) \cdot (1 + \delta \cdot
  \sin(\omega_{\mathrm{Ca}} t))$.
\item \textit{Phase shift:}
  $\Delta\phi_{\mathrm{Ca}} = \delta \cdot
  \sin(\omega_{\mathrm{Ca}} t_{\mathrm{obs}})$.
\item \textit{Effect on ray:} oscillatory modulation of circuit
  current $\Jray$.
\end{enumerate}
\end{definition}

\subsubsection{Circadian Oscillators ($\sim\!10^{-5}$\,Hz):
Slow Parametric Drift}

The circadian clock, with a period of approximately 24 hours,
modulates global cellular parameters \citep{reppert2002}.

\begin{definition}[Circadian Interaction]
\label{def:circadian-interaction}
\begin{enumerate}[label=(\alph*),leftmargin=*]
\item \textit{Partition signature:} all parameters drift slowly with
  circadian phase $\Phi_R \in [0, 2\pi)$.
\item \textit{Interaction:} parametric drift of the form
  $\alpha(t) = \alpha_0(1 + \epsilon_R \cos(\omega_R t))$ for every
  interaction parameter $\alpha$.
\item \textit{Phase shift:} $\Delta\phi_R \approx \epsilon_R \ll 1$
  per ray traversal.
\item \textit{Effect on ray:} negligible on the time scale of a single
  ray traversal, but detectable as a systematic drift over many
  observations spanning hours.
\end{enumerate}
\end{definition}

% ── Figure 7 proposal ──────────────────────────────────────────────
% FIGURE 7: The eight oscillator classes as a frequency spectrum.
% Horizontal axis: log frequency from 10^{-5} to 10^{14} Hz.
% Each class shown as a colored band. Above each band: the ray
% interaction type. Below each band: an example biological process.


% ═════════════════════════════════════════════════════════════════════
\section{Phase Accumulation and Reaction History}
\label{sec:phase-accumulation}

\begin{theorem}[Phase Accumulation]
\label{thm:phase-accumulation}
Consider a ray traversing a region where bioreaction
$\Gamma_1 \oplus P(\omega) \to \Gamma_2$ occurs at oscillator
frequency $\omega$ with partition lag $\taup$. The ray accumulates
a phase shift:
\begin{equation}
\Delta\phi = \omega \cdot \taup.
\label{eq:phase-shift}
\end{equation}
The total accumulated phase along the entire ray path is:
\begin{equation}
\Phi_{\mathrm{total}} = \int_{\mathrm{path}} \sum_{k=1}^{8}
  \omega_k(\mathbf{r}(s)) \cdot \taup^{(k)}(\mathbf{r}(s)) \, ds,
\label{eq:total-phase}
\end{equation}
where the sum runs over the eight oscillator classes and the integral
runs over the ray path. This accumulated phase encodes the complete
reaction history.
\end{theorem}

\begin{proof}
Each bioreaction at position $\mathbf{r}$ along the ray path imparts
a phase shift $\Delta\phi = \omega \cdot \taup$ to the ray. The
product $\omega \cdot \taup$ is the number of radians of oscillation
that occur during one partition transition --- it measures the
``depth'' of the interaction.

For protein oscillators: $\omega_P \sim 10^{13}$\,Hz, $\taup^{(P)}
\sim 10^{-13}$\,s, so $\Delta\phi_P \sim 1$\,rad.

For enzyme oscillators: $\omega_E \sim 10^{8}$\,Hz, $\taup^{(E)}
\sim 10^{-6}$\,s, so $\Delta\phi_E \sim 100$\,rad (enzymes are
``deep'' scatterers).

For channel oscillators: $\omega_C \sim 10^{4}$\,Hz, $\taup^{(C)}
\sim 10^{-3}$\,s, so $\Delta\phi_C \sim 10$\,rad.

The total accumulated phase is the line integral of the local phase
accumulation rate $\sum_k \omega_k \taup^{(k)}$ along the path. Since
the integrand depends on position through the partition state
$\Sigma(\mathbf{r})$, the accumulated phase $\Phi_{\mathrm{total}}$ is
a functional of the path and encodes what oscillators the ray
encountered, in what order, and for how long --- i.e., the complete
reaction history along the path.
\end{proof}

\begin{corollary}[Phase Encodes Partition History]
\label{cor:phase-history}
Two rays that traverse the same sequence of partition states
accumulate the same phase. Two rays that traverse different partition
sequences accumulate different phases. The accumulated phase is a
unique fingerprint of the partition history.
\end{corollary}


% ═════════════════════════════════════════════════════════════════════
\section{Coherence from Interference}
\label{sec:coherence}

\subsection{Multi-Ray Interference}

\begin{definition}[Ray Ensemble]
\label{def:ray-ensemble}
Launch $N$ rays through the cellular volume along different paths
$\{\gamma_1, \gamma_2, \ldots, \gamma_N\}$. Each ray $k$ accumulates
a total phase $\Phi_k$ from the oscillators it encounters along path
$\gamma_k$.
\end{definition}

\begin{definition}[Interference Visibility]
\label{def:visibility}
The interference visibility of the ray ensemble is:
\begin{equation}
\Vcell = \left| \langle e^{i\Phi_k} \rangle \right|
  = \left| \frac{1}{N} \sum_{k=1}^{N} e^{i\Phi_k} \right|.
\label{eq:visibility}
\end{equation}
$\Vcell \in [0, 1]$, with $\Vcell = 1$ when all rays accumulate
the same phase (perfect coherence) and $\Vcell = 0$ when phases are
uniformly distributed (complete decoherence).
\end{definition}

\subsection{Coherence as Health Diagnostic}

\begin{theorem}[Coherence--Health Identity]
\label{thm:coherence-health}
The interference visibility $\Vcell$ is identical to the cellular
coherence index $\etacell$:
\begin{equation}
\boxed{
\Vcell = \etacell = \frac{1}{W} \sum_{i,j} w_{ij} \cdot
  \frac{\Pi_{ij,\mathrm{obs}} - \Pi_{ij,\mathrm{deg}}}
  {\Pi_{ij,\mathrm{opt}} - \Pi_{ij,\mathrm{deg}}},
}
\label{eq:coherence-health}
\end{equation}
where $W = \sum_{i,j} w_{ij}$, the sum runs over all pairs $(i,j)$
of oscillator classes, $w_{ij}$ is the coupling weight between classes
$i$ and $j$, $\Pi_{ij,\mathrm{obs}}$ is the observed coherence between
oscillators $i$ and $j$, and $\Pi_{ij,\mathrm{opt}}$ and
$\Pi_{ij,\mathrm{deg}}$ are the optimal (healthy) and degraded
(pathological) coherences, respectively.
\end{theorem}

\begin{proof}
The proof proceeds by showing that both sides measure the same
quantity: the degree of synchronization among the cellular oscillators.

\textit{Left side (ray interference):} When the eight oscillator
classes are well-synchronized (healthy cell), the phase shifts they
impart to different rays are correlated. Specifically, for two rays
$k$ and $\ell$ passing through the same oscillator at different
positions, the phase difference $\Phi_k - \Phi_\ell$ is small if the
oscillator is coherent (same frequency and fixed phase relationship
across the cell) and large if the oscillator is decoherent (variable
frequency or random phase). The ensemble average
$|\langle e^{i\Phi_k} \rangle|$ is therefore large when oscillators
are coherent and small when they are decoherent.

To make this precise, decompose the total phase by oscillator class:
\begin{equation}
\Phi_k = \sum_{i=1}^{8} \Phi_k^{(i)},
\end{equation}
where $\Phi_k^{(i)}$ is the phase accumulated from class-$i$
oscillators along path $k$. Then:
\begin{align}
\langle e^{i\Phi_k} \rangle &= \left\langle
  \prod_{i=1}^{8} e^{i\Phi_k^{(i)}} \right\rangle \notag \\
  &= \prod_{i=1}^{8} \langle e^{i\Phi_k^{(i)}} \rangle
  + \text{cross-terms}.
\label{eq:phase-decomposition}
\end{align}

The cross-terms involve products
$\langle e^{i\Phi_k^{(i)}} e^{i\Phi_k^{(j)}} \rangle$ for $i \neq j$,
which are the inter-class coherences. Defining:
\begin{equation}
\Pi_{ij,\mathrm{obs}} = \left| \left\langle e^{i(\Phi_k^{(i)} +
  \Phi_k^{(j)})} \right\rangle \right|,
\end{equation}
we can write:
\begin{equation}
\Vcell = \frac{1}{W} \sum_{i \leq j} w_{ij} \cdot \Pi_{ij,\mathrm{obs}},
\end{equation}
where $w_{ij}$ weights the contribution of each pair.

\textit{Right side (coherence index):} The normalization
$(\Pi_{ij,\mathrm{obs}} - \Pi_{ij,\mathrm{deg}}) /
(\Pi_{ij,\mathrm{opt}} - \Pi_{ij,\mathrm{deg}})$ maps each pair
coherence to $[0,1]$, with $0$ for the completely decoherent
(pathological) state and $1$ for the optimally coherent (healthy)
state.

The identification $\Vcell = \etacell$ follows when we set:
\begin{equation}
\Pi_{ij,\mathrm{deg}} = 0 \quad \text{and} \quad
\Pi_{ij,\mathrm{opt}} = 1,
\end{equation}
corresponding to zero coherence (uniformly random phases) and perfect
coherence (identical phases), respectively. With this normalization,
$\etacell = (1/W) \sum w_{ij} \Pi_{ij,\mathrm{obs}} = \Vcell$.

For biological systems where the degenerate and optimal coherences
are not exactly $0$ and $1$, the normalization adjusts accordingly,
but the identity between ray visibility and oscillator coherence
is maintained.
\end{proof}

\begin{corollary}[Diagnostic Interpretation]
\label{cor:diagnostic}
\begin{enumerate}[label=(\roman*),leftmargin=*]
\item $\etacell \approx 1$: all oscillator classes are synchronized.
  The cell is healthy. Rays interfere constructively. The visibility
  is high.
\item $\etacell \approx 0$: oscillators are desynchronized. The cell
  is pathological. Rays interfere destructively. The visibility is
  low.
\item Intermediate values: partial coherence. Specific decoherent
  oscillator classes can be identified from the pair coherences
  $\Pi_{ij}$, enabling targeted diagnosis.
\end{enumerate}
\end{corollary}

% ── Figure 8 proposal ──────────────────────────────────────────────
% FIGURE 8: Two panels. Left: healthy cell — rays shown with arrows,
% phase shown as color (all similar), interference pattern below shows
% high visibility. Right: diseased cell — rays with random phase
% colors, interference pattern shows low visibility. Center bar:
% eta_cell scale from 0 (red) to 1 (green).


% ═════════════════════════════════════════════════════════════════════
%             PART IV: HOLOGRAPHIC 3D RECONSTRUCTION
% ═════════════════════════════════════════════════════════════════════

\section{The Dual-Pixel Hologram}
\label{sec:dual-pixel}

To reconstruct the three-dimensional cellular volume from a
two-dimensional observation, we require both amplitude and phase
information --- i.e., a holographic record \citep{gabor1948}.

\subsection{Visible and Invisible Pixels}

\begin{definition}[Dual-Pixel Structure]
\label{def:dual-pixel}
Each observation pixel consists of two components:
\begin{enumerate}[label=(\roman*),leftmargin=*]
\item The \emph{visible pixel}: the standard optical signal recording
  amplitude $|U(x,y)|$ (intensity, color, polarization).
\item The \emph{invisible pixel}: the oxygen-state signal recording
  phase via the emission timing
  $\tauem = 1/A_{ki}$, where $A_{ki}$ is the Einstein spontaneous
  emission coefficient \citep{einstein1917}. The O$_2$ molecule at
  position $(x,y)$ encodes its local chemical environment as a ternary
  state (triplet ground state ${}^3\Sigma_g^-$, singlet excited
  ${}^1\Delta_g$, or singlet excited ${}^1\Sigma_g^+$), and the
  emission timing from each state encodes the phase of the local
  oscillatory environment.
\end{enumerate}
\end{definition}

\begin{theorem}[Holographic Completeness]
\label{thm:holographic}
The dual-pixel structure (visible + invisible) constitutes a complete
holographic record: together, they encode both amplitude $|U|$ and
phase $\arg(U)$ of the wavefield at the observation plane.
\end{theorem}

\begin{proof}
A hologram is defined as a recording of both the amplitude and phase
of a wavefield \citep{gabor1948,goodman2005}. The visible pixel
records $|U(x,y)|^2 \propto I(x,y)$, from which $|U|$ can be
extracted.

The invisible pixel records phase through the emission timing. The
Einstein $A$ coefficient depends on the local electromagnetic
environment (Purcell effect), which is modulated by the oscillatory
state of the medium. Specifically:
\begin{equation}
\tauem(\mathbf{r}) = \frac{1}{A_{ki}(\mathbf{r})}
  = \tau_0 \cdot \left(1 + \beta \sum_j
  \frac{\omega_j(\mathbf{r})}{\omega_0}\right)^{-1},
\label{eq:emission-timing}
\end{equation}
where $\tau_0$ is the free-space emission time, $\beta$ is the
coupling strength, and the sum runs over local oscillators. The
emission timing is therefore a function of the local phase environment.

By measuring $\tauem$ at each pixel position $(x,y)$, we obtain a
phase map:
\begin{equation}
\phi(x,y) = \omega_0 \cdot \tauem(x,y) \pmod{2\pi}.
\label{eq:phase-map}
\end{equation}

The complex field is then:
\begin{equation}
U(x,y) = |U(x,y)| \cdot e^{i\phi(x,y)},
\label{eq:complex-field}
\end{equation}
which is the complete holographic record.
\end{proof}


% ═════════════════════════════════════════════════════════════════════
\section{Angular Spectrum Propagation}
\label{sec:angular-spectrum}

Given the two-dimensional holographic field $U(x,y,0)$ at the
observation plane $z=0$, we reconstruct the three-dimensional field by
angular spectrum back-propagation \citep{goodman2005,angularspectrum1968}.

\begin{theorem}[3D Reconstruction by Angular Spectrum]
\label{thm:angular-spectrum}
The three-dimensional field at any depth $z$ is:
\begin{equation}
U(x,y,z) = \mathcal{F}^{-1}\!\left\{
  \hat{U}(k_x, k_y) \cdot e^{ik_z z}
\right\},
\label{eq:angular-spectrum}
\end{equation}
where $\hat{U}(k_x, k_y) = \mathcal{F}\{U(x,y,0)\}$ is the 2D
Fourier transform of the observation-plane field, and:
\begin{equation}
k_z = \sqrt{k_0^2 - k_x^2 - k_y^2},
\label{eq:kz}
\end{equation}
with $k_0 = 2\pi/\lambda$ the free-space wavenumber. For
$k_x^2 + k_y^2 > k_0^2$, $k_z$ is imaginary (evanescent waves), and
these components decay exponentially with $z$.
\end{theorem}

\begin{proof}
This is the standard angular spectrum representation of diffraction
\citep{bornwolf1999,goodman2005}. Any monochromatic field satisfying
the Helmholtz equation $(\nabla^2 + k_0^2)U = 0$ can be decomposed
into plane-wave components $e^{i(k_x x + k_y y + k_z z)}$. The
field at $z=0$ determines the angular spectrum $\hat{U}(k_x, k_y)$,
and propagation to depth $z$ multiplies each component by the
propagation kernel $e^{ik_z z}$. Back-propagation (reconstruction from
$z=0$ data to $z > 0$) is achieved by using the conjugate kernel
$e^{-ik_z z}$ or equivalently by using the forward kernel with $z > 0$
for the reverse direction.
\end{proof}

\begin{remark}[GPU Implementation]
Each depth slice $z_j$ requires one 2D FFT (to compute
$\hat{U}(k_x, k_y)$), one element-wise multiply by $e^{ik_z z_j}$,
and one 2D inverse FFT. For a $256 \times 256$ observation plane and
$256$ depth slices, this is $256$ FFT pairs, each of size
$256 \times 256$. On a modern GPU, each FFT pair takes
$\sim\!0.1$\,ms, so the total reconstruction time is
$\sim\!25.6$\,ms ($\approx 39$ depth slices per millisecond). This is
well within real-time constraints.
\end{remark}


% ═════════════════════════════════════════════════════════════════════
\section{Volumetric Partition State}
\label{sec:volumetric-state}

\begin{theorem}[Partition State from Reconstructed Field]
\label{thm:partition-from-field}
The reconstructed three-dimensional field $U(x,y,z)$ encodes the
complete volumetric partition signature
$\Sigma(x,y,z) = (n, \ell, m, s)$ at every three-dimensional
position through the relations:
\begin{align}
n(x,y,z) &= n_{\max} \cdot |U(x,y,z)|^2 / |U|^2_{\max},
  \label{eq:n-from-field} \\
\ell(x,y,z) &= \lfloor n \cdot |\nabla_\perp \arg(U)|
  / |\nabla_\perp \arg(U)|_{\max} \rfloor,
  \label{eq:ell-from-field} \\
m(x,y,z) &= \lfloor (2\ell+1) \cdot \arg(U)/(2\pi) \rfloor - \ell,
  \label{eq:m-from-field} \\
s(x,y,z) &= \sgn(\partial_z |U|^2).
  \label{eq:s-from-field}
\end{align}
\end{theorem}

\begin{proof}
\textit{Principal depth $n$:} The intensity $|U|^2$ at each voxel
measures the local ``brightness,'' which in the partition framework
corresponds to categorical density --- higher intensity means more
categorical states are populated, corresponding to higher $n$.
Normalization by $|U|^2_{\max}$ maps to $[0, n_{\max}]$.

\textit{Angular index $\ell$:} The transverse phase gradient
$|\nabla_\perp \arg(U)|$ measures the local wavefront curvature,
which corresponds to the angular complexity of the partition state.
Higher curvature means higher $\ell$.

\textit{Orientation index $m$:} The phase angle $\arg(U)$ itself
determines the orientation within the $\ell$ subpartition. Quantization
to the $2\ell+1$ allowed values and centering at $-\ell$ gives the
standard range $m \in \{-\ell, \ldots, \ell\}$.

\textit{Parity $s$:} The sign of the axial intensity gradient
$\partial_z |U|^2$ determines whether the local environment is
``absorbing'' ($|U|^2$ decreasing along $z$, $s = -1/2$) or
``emitting/scattering'' ($|U|^2$ increasing, $s = +1/2$).
\end{proof}

% ── Figure 9 proposal ──────────────────────────────────────────────
% FIGURE 9: The holographic reconstruction pipeline. (a) 2D
% observation plane with amplitude and phase. (b) Angular spectrum
% in Fourier space. (c) Propagated to depth z_1. (d) Full 3D volume
% with partition coordinates color-coded.


% ═════════════════════════════════════════════════════════════════════
%             PART V: THE GPU RAY TRACER ARCHITECTURE
% ═════════════════════════════════════════════════════════════════════

\section{Volumetric Data Structure}
\label{sec:volumetric-data}

\begin{definition}[Cellular Volume Texture]
\label{def:volume-texture}
The three-dimensional cellular volume is stored as a 3D texture
(or equivalently, a stack of layered 2D textures for WebGL~2.0
compatibility \citep{khronos2024}) with each voxel storing the
partition coordinates:
\begin{equation}
\text{Voxel}(i,j,k) = (n, \ell, m, s) \in \R^4
  \quad \text{stored as RGBA32F}.
\label{eq:voxel}
\end{equation}
\end{definition}

\begin{proposition}[Memory Requirements]
\label{prop:memory}
The memory footprint for common volume resolutions:
\begin{center}
\small
\begin{tabular}{@{}lcc@{}}
\toprule
\textbf{Resolution} & \textbf{Voxels} & \textbf{Memory (MB)} \\
\midrule
$64^3$ & 262\,144 & 4 \\
$128^3$ & 2\,097\,152 & 32 \\
$256^3$ & 16\,777\,216 & 256 \\
\bottomrule
\end{tabular}
\end{center}
At RGBA32F (16 bytes per voxel): $V^3 \times 16$ bytes. A $128^3$
volume at 32\,MB fits comfortably in any integrated GPU with $\geq
128$\,MB shared memory (all modern integrated GPUs have at least
this) \citep{owens2008}.
\end{proposition}


% ═════════════════════════════════════════════════════════════════════
\section{Ray March Shader}
\label{sec:ray-march}

\begin{algorithm}[H]
\caption{Triple-Observation Ray March}
\label{alg:ray-march}
\begin{algorithmic}[1]
\Require Ray origin $\mathbf{o}$, ray direction $\hat{\mathbf{d}}$,
  step size $\Delta s$, max steps $N_{\max}$, transmittance threshold
  $T_{\min}$
\Ensure Accumulated color $\mathbf{C}$, retention $t_R$, current $J$,
  phase $\Phi$
\State $\mathbf{C} \gets (0,0,0)$; $T \gets 1$; $t_R \gets 0$;
  $J \gets 0$; $\Phi \gets 0$
\State $\Psi_{\mathrm{ray}} \gets \Psi_0$ \Comment{Initial ray
  S-entropy state}
\State $\varphi_{\mathrm{ray}} \gets \varphi_0$ \Comment{Initial ray
  reference potential}
\For{$i = 0$ \textbf{to} $N_{\max}$}
  \State $\mathbf{r} \gets \mathbf{o} + i \cdot \Delta s \cdot
    \hat{\mathbf{d}}$ \Comment{Current position}
  \If{$\mathbf{r} \notin \text{volume}$} \textbf{break}
  \EndIf
  \State $\Sigma \gets \text{texture3D}(\mathbf{r})$
    \Comment{Read $(n,\ell,m,s)$}
  \State \Comment{\textbf{--- Optical observation ---}}
  \State $\muabs \gets \alpha \cdot (n/n_{\max}) \cdot F(\ell,m,s)$
  \State $T \gets T \cdot \exp(-\muabs \cdot \Delta s)$
  \State $\mathbf{C} \gets \mathbf{C} + T \cdot \muabs \cdot
    \Delta s \cdot \mathbf{c}_{\mathrm{emit}}(\Sigma)$
  \State \Comment{\textbf{--- Chromatographic observation ---}}
  \State $d \gets \dS(\Psi_{\mathrm{ray}}, \Sigma)$
  \State $t_R \gets t_R + \tau \cdot d \cdot \Delta s$
  \State $\Psi_{\mathrm{ray}} \gets \Psi_{\mathrm{ray}} +
    \lambda_{\mathrm{adapt}} \cdot (\Sigma_S - \Psi_{\mathrm{ray}})
    \cdot \Delta s$ \Comment{Ray adapts to medium}
  \State \Comment{\textbf{--- Circuit observation ---}}
  \State $\varphi_{\mathrm{local}} \gets -\kb T \ln 2 \cdot
    H_{\mathrm{cat}}(\Sigma) + c$
  \State $G_{\mathrm{local}} \gets g_0 \cdot (n/n_{\max}) \cdot F$
  \State $J \gets J + G_{\mathrm{local}} \cdot
    (\varphi_{\mathrm{local}} - \varphi_{\mathrm{ray}}) \cdot \Delta s$
  \State \Comment{\textbf{--- Phase accumulation ---}}
  \State $\omega_{\mathrm{local}} \gets \sum_k \omega_k(\Sigma)$
    \Comment{Sum over oscillator classes}
  \State $\Phi \gets \Phi + \omega_{\mathrm{local}} \cdot \taup(\Sigma)
    \cdot \Delta s$
  \State \Comment{\textbf{--- Early termination ---}}
  \If{$T < T_{\min}$} \textbf{break}
  \EndIf
\EndFor
\State \Return $(\mathbf{C}, t_R, J, \Phi)$
\end{algorithmic}
\end{algorithm}

The following GLSL pseudocode implements Algorithm~\ref{alg:ray-march}:

\begin{lstlisting}[style=glsl,caption={Triple-observation ray march
  fragment shader (GLSL pseudocode)},
  label={lst:ray-march}]
// Triple-Observation Ray March Shader
// Computes optical, chromatographic, and
// circuit observations simultaneously

uniform sampler3D u_volume; // 3D partition texture
uniform vec3 u_ray_origin;
uniform vec3 u_ray_dir;
uniform float u_step_size;
uniform int u_max_steps;
uniform float u_alpha_max;   // max absorption
uniform float u_tau;         // column time constant
uniform float u_g0;          // ref conductance
uniform float u_n_max;       // max principal depth
uniform float u_T_min;       // early termination

out vec4 fragColor;  // .rgb = color, .a = transmit
out vec4 fragData1;  // .r = retention, .g = current
                     // .b = phase, .a = unused

void main() {
  vec3 pos = u_ray_origin;
  vec3 dir = normalize(u_ray_dir);
  float ds = u_step_size;

  // Accumulated quantities
  vec3 color = vec3(0.0);
  float transmittance = 1.0;
  float retention = 0.0;
  float current = 0.0;
  float phase = 0.0;

  // Ray state
  vec3 psi_ray = vec3(0.33, 0.33, 0.34);
  float phi_ray = 0.0;

  for (int i = 0; i < u_max_steps; i++) {
    pos += dir * ds;

    // Bounds check
    if (any(lessThan(pos, vec3(0.0))) ||
        any(greaterThan(pos, vec3(1.0))))
      break;

    // Read partition state from 3D texture
    vec4 sigma = texture(u_volume, pos);
    float n = sigma.r;   // principal depth
    float ell = sigma.g; // angular index
    float m = sigma.b;   // orientation
    float s = sigma.a;   // parity

    // --- Optical observation ---
    float F = (1.0 - ell/max(n, 0.01))
            * cos(3.14159*m/(2.0*ell+1.0))
            * cos(3.14159*m/(2.0*ell+1.0))
            * step(0.0, s);
    float mu_a = u_alpha_max * (n/u_n_max) * F;
    transmittance *= exp(-mu_a * ds);
    vec3 emit = vec3(n/u_n_max,
                     ell/max(n,0.01),
                     abs(m)/(ell+0.01));
    color += transmittance * mu_a * ds * emit;

    // --- Chromatographic observation ---
    vec3 sigma_S = vec3(
      (1.0 - ell/max(n,0.01)),
      ell/max(n,0.01),
      abs(s));
    float d_S = distance(psi_ray, sigma_S);
    retention += u_tau * d_S * ds;
    psi_ray = mix(psi_ray, sigma_S, 0.01*ds);

    // --- Circuit observation ---
    float H_cat = -n/u_n_max
                  * log2(max(n/u_n_max, 0.001));
    float phi_local = -1.38e-23 * 310.0
                      * 0.693 * H_cat;
    float G_local = u_g0 * (n/u_n_max) * F;
    current += G_local
             * (phi_local - phi_ray) * ds;

    // --- Phase accumulation ---
    float omega = n * 1e13 / u_n_max;
    float tau_p = 1.0 / max(omega, 1.0);
    phase += omega * tau_p * ds;

    // --- Early termination ---
    if (transmittance < u_T_min) break;
  }

  fragColor = vec4(color, transmittance);
  fragData1 = vec4(retention, current,
                   phase, 0.0);
}
\end{lstlisting}

\begin{remark}[Computational Cost]
Each ray march step performs: one texture fetch (the dominant cost on
integrated GPUs), approximately 30 arithmetic operations (multiply,
add, exp, log, distance), and two accumulations. For a $512 \times 512$
screen with 256 steps per ray, the total operation count is
$512^2 \times 256 \times 30 \approx 2 \times 10^9$ FLOP. At 1\,TFLOPS
(typical integrated GPU), this takes $\sim\!2$\,ms per frame --- well
within real-time constraints.
\end{remark}


% ═════════════════════════════════════════════════════════════════════
\section{Multi-Ray Interference Shader}
\label{sec:multi-ray}

\begin{algorithm}[H]
\caption{Multi-Ray Interference for Coherence Extraction}
\label{alg:multi-ray}
\begin{algorithmic}[1]
\Require Ray results $\{(\mathbf{C}_k, t_{R,k}, J_k,
  \Phi_k)\}_{k=1}^N$ from $N$ ray directions
\Ensure 3D morphology, metabolic flow map, circuit state, health
  diagnostic
\State \Comment{\textbf{--- 3D Morphology (from amplitude) ---}}
\State $\text{Volume}(x,y,z) \gets \frac{1}{N}\sum_k
  |\mathbf{C}_k(x,y,z)|$ \Comment{Average color magnitude}
\State \Comment{\textbf{--- Metabolic Flow Map (from retention) ---}}
\State $\text{FlowMap}(x,y,z) \gets \nabla \langle t_{R,k}
  \rangle$ \Comment{Retention gradient}
\State \Comment{\textbf{--- Circuit State (from current) ---}}
\State $\text{CircuitMap}(x,y,z) \gets \langle J_k \rangle$
  \Comment{Average accumulated current}
\State \Comment{\textbf{--- Health Diagnostic (from interference) ---}}
\State $\Vcell \gets \left|\frac{1}{N}\sum_{k=1}^N e^{i\Phi_k}
  \right|$ \Comment{Interference visibility}
\State $\etacell \gets \Vcell$ \Comment{Coherence index}
\State \Return $(\text{Volume}, \text{FlowMap}, \text{CircuitMap},
  \etacell)$
\end{algorithmic}
\end{algorithm}

\begin{lstlisting}[style=glsl,caption={Multi-ray interference shader
  (GLSL pseudocode)},
  label={lst:multi-ray}]
// Multi-Ray Interference Shader
// Combines results from N ray directions

uniform sampler2D u_phase_maps[8];
  // Phase maps from 8 ray directions
uniform int u_num_rays;

out vec4 fragCoherence;

void main() {
  vec2 uv = gl_FragCoord.xy
           / vec2(textureSize(u_phase_maps[0],0));

  // Compute interference visibility
  float sum_cos = 0.0;
  float sum_sin = 0.0;

  for (int k = 0; k < u_num_rays; k++) {
    float phi_k = texture(u_phase_maps[k], uv).b;
    sum_cos += cos(phi_k);
    sum_sin += sin(phi_k);
  }

  float N = float(u_num_rays);
  float V_cell = sqrt(sum_cos*sum_cos
                    + sum_sin*sum_sin) / N;

  // V_cell IS the coherence index eta_cell
  float eta_cell = V_cell;

  // Extract per-oscillator-class coherences
  // (requires decomposing phase by frequency)
  // ...omitted for brevity...

  fragCoherence = vec4(eta_cell,
    V_cell, 0.0, 1.0);
}
\end{lstlisting}


% ═════════════════════════════════════════════════════════════════════
\section{Complete Pipeline}
\label{sec:pipeline}

The complete GPU pipeline consists of five passes:

\begin{definition}[Five-Pass Pipeline]
\label{def:pipeline}
\begin{enumerate}[label=\textbf{Pass~\arabic*:},leftmargin=*]
\item[\textbf{Pass 0:}] \textbf{Dual-Pixel Analysis.}
  Input: 2D microscopy image. The fragment shader analyzes each pixel
  to extract amplitude $|U(x,y)|$ from the visible channel and phase
  $\phi(x,y)$ from the emission timing of the invisible (O$_2$)
  channel. Output: complex field texture $U(x,y) = |U| e^{i\phi}$.
  Cost: one draw call, $\sim\!0.1$\,ms.

\item[\textbf{Pass 1:}] \textbf{Holographic Back-Propagation.}
  Input: $U(x,y,0)$ from Pass~0. For each depth slice $z_j$: compute
  2D FFT, multiply by $e^{ik_z z_j}$, compute inverse FFT. Store
  result in 3D texture. Output: volumetric partition state
  $\Sigma(x,y,z)$ via Theorem~\ref{thm:partition-from-field}. Cost:
  $\sim\!25$\,ms for $256^3$.

\item[\textbf{Pass 2:}] \textbf{Triple-Observation Ray March.}
  Input: 3D volume texture from Pass~1. For each pixel on screen:
  march a ray through the volume using
  Algorithm~\ref{alg:ray-march}. Output: per-pixel color,
  retention, current, phase (stored in two render targets).
  Repeat for $N$ ray directions ($N = 8$ gives good coverage). Cost:
  $\sim\!2$\,ms per direction, $\sim\!16$\,ms total.

\item[\textbf{Pass 3:}] \textbf{Multi-Ray Interference.}
  Input: $N$ phase maps from Pass~2. Compute interference visibility
  per pixel using Algorithm~\ref{alg:multi-ray}. Output: coherence
  map $\etacell(x,y)$. Cost: $\sim\!0.5$\,ms.

\item[\textbf{Pass 4:}] \textbf{Scalar Readback.}
  Read back scalar diagnostics: mean coherence index
  $\bar{\eta}_{\mathrm{cell}}$, total metabolic flux, circuit state
  summary. This uses GPU readback (one small texture). Cost:
  $\sim\!1$\,ms.
\end{enumerate}
\end{definition}

\begin{proposition}[Total Pipeline Latency]
\label{prop:latency}
The total pipeline latency is:
\begin{equation}
T_{\mathrm{total}} = T_{\mathrm{P0}} + T_{\mathrm{P1}}
  + N \cdot T_{\mathrm{P2}} + T_{\mathrm{P3}} + T_{\mathrm{P4}}
  \approx 0.1 + 25 + 16 + 0.5 + 1 = 42.6\,\text{ms},
\label{eq:total-latency}
\end{equation}
corresponding to $\sim\!23$ frames per second. This is real-time on an
integrated GPU with $\sim\!1$\,TFLOPS compute and $\sim\!32$\,MB
working set.
\end{proposition}

\begin{proposition}[Memory Budget]
\label{prop:memory-budget}
\begin{center}
\small
\begin{tabular}{@{}lc@{}}
\toprule
\textbf{Resource} & \textbf{Memory (MB)} \\
\midrule
3D volume texture ($128^3$, RGBA32F) & 32 \\
2D observation texture ($512^2$, RGBA32F) & 4 \\
$N=8$ phase map textures ($512^2$, R32F) & 8 \\
Shader programs and uniforms & $<\!1$ \\
\midrule
\textbf{Total} & $\sim\!45$ \\
\bottomrule
\end{tabular}
\end{center}
This fits within the shared memory of any integrated GPU manufactured
after 2015.
\end{proposition}

% ── Figure 10 proposal ─────────────────────────────────────────────
% FIGURE 10: The five-pass pipeline as a flow diagram. Pass 0 → Pass
% 1 → Pass 2 (×N) → Pass 3 → Pass 4. Each box shows input/output
% textures and timing. Total: ~43 ms.


% ═════════════════════════════════════════════════════════════════════
%             PART VI: FLUID DYNAMICS FROM RAY TRACING
% ═════════════════════════════════════════════════════════════════════

\section{Cytoplasmic Flow as Chromatographic Transport}
\label{sec:cytoplasmic-flow}

The cytoplasm is not static: it has a velocity field
$\mathbf{v}(\mathbf{r})$ driven by molecular motors, osmotic
pressure, and active transport. We now show that the ray tracer
implicitly measures this flow.

\subsection{Doppler-Shifted Retention}

\begin{theorem}[Flow-Dependent Retention]
\label{thm:flow-retention}
A ray traversing a moving medium with velocity field
$\mathbf{v}(\mathbf{r})$ experiences a Doppler-shifted categorical
distance:
\begin{equation}
\dS^{\mathrm{eff}} = \dS \cdot \left(1 +
  \frac{\mathbf{v} \cdot \hat{\mathbf{r}}}{\cS}\right),
\label{eq:doppler-retention}
\end{equation}
where $\hat{\mathbf{r}}$ is the ray direction and $\cS$ is the
categorical signal velocity.
\end{theorem}

\begin{proof}
The categorical distance measures the ``interaction strength'' between
the ray and the local medium. When the medium moves with velocity
$\mathbf{v}$, the effective interaction time depends on the relative
velocity between the ray and the medium.

In a frame co-moving with the medium, the ray has velocity
$c - \mathbf{v} \cdot \hat{\mathbf{r}}$ (where $c$ is the ray
propagation speed). The interaction time is inversely proportional to
this speed, so the effective retention scales as:
\begin{equation}
\dS^{\mathrm{eff}} = \dS \cdot \frac{c}{c -
  \mathbf{v} \cdot \hat{\mathbf{r}}}.
\end{equation}
For $|\mathbf{v}| \ll c$, this linearizes to:
\begin{equation}
\dS^{\mathrm{eff}} \approx \dS \cdot \left(1 +
  \frac{\mathbf{v} \cdot \hat{\mathbf{r}}}{c}\right).
\end{equation}
We identify $c$ with the categorical signal velocity $\cS$, which is
the speed at which partition state changes propagate through the
medium. For cellular systems, $\cS$ is determined by the speed of
propagation of conformational changes along protein networks and
membrane signaling, estimated at $\cS \sim 10^3$\,m/s from the
propagation speed of allosteric transitions \citep{groel1997}.
\end{proof}

\begin{corollary}[Flow Velocity from Retention Anisotropy]
\label{cor:flow-velocity}
By launching rays in multiple directions and comparing their retention
times, the local flow velocity can be extracted:
\begin{equation}
\mathbf{v}(\mathbf{r}) = \cS \cdot \frac{\dS^{\mathrm{eff}}
  (\hat{\mathbf{r}}_+) - \dS^{\mathrm{eff}}(\hat{\mathbf{r}}_-)}
  {\dS^{\mathrm{eff}}(\hat{\mathbf{r}}_+) +
  \dS^{\mathrm{eff}}(\hat{\mathbf{r}}_-)},
\label{eq:flow-from-retention}
\end{equation}
where $\hat{\mathbf{r}}_+$ and $\hat{\mathbf{r}}_-$ are opposing ray
directions. This is the chromatographic analogue of Doppler
velocimetry.
\end{corollary}


% ═════════════════════════════════════════════════════════════════════
\section{Navier--Stokes from Kirchhoff}
\label{sec:navier-stokes}

\begin{theorem}[Stokes Equations from Circuit Laws]
\label{thm:stokes-from-kirchhoff}
The cellular circuit's Kirchhoff laws, extended to a continuous spatial
field, yield the Stokes equations appropriate for the cellular interior.
\end{theorem}

\begin{proof}
\textit{Step 1: Kirchhoff's Current Law $\to$ Continuity Equation.}

Kirchhoff's current law at each node (Eq.~\eqref{eq:kcl}) states that
the sum of currents (reaction fluxes) into each node equals the source
term. In the continuous limit, where nodes become spatial positions and
currents become flux density vectors:
\begin{equation}
\sum_j G_{ij}(\varphi_i - \varphi_j) \to
  \nabla \cdot [G(\mathbf{r}) \nabla\varphi(\mathbf{r})]
  = Q(\mathbf{r}).
\label{eq:continuous-kcl}
\end{equation}

Identifying the chemical potential gradient with the pressure gradient
$\nabla p$ and the flux with the velocity $\mathbf{v} = -G \nabla
\varphi / \rho$, this becomes:
\begin{equation}
\frac{\partial \rho}{\partial t} +
  \nabla \cdot (\rho \mathbf{v}) = Q(\mathbf{r}),
\label{eq:continuity}
\end{equation}
which is the continuity equation \citep{navierstokes}.

\textit{Step 2: Kirchhoff's Voltage Law $\to$ Pressure Constraint.}

Kirchhoff's voltage law (Eq.~\eqref{eq:kvl}) states that around any
closed loop, the sum of potential drops is zero. In the continuous
limit, this becomes the irrotationality condition on the potential
gradient:
\begin{equation}
\oint \nabla\varphi \cdot d\mathbf{l} = 0
  \quad \Longrightarrow \quad
  \nabla \times (\nabla\varphi) = 0,
\label{eq:continuous-kvl}
\end{equation}
which is automatically satisfied for any scalar potential. This
constrains the pressure field to be a gradient field.

\textit{Step 3: Low Reynolds Number.}

The Reynolds number for cytoplasmic flow is
\citep{purcell1977,bergetal2002}:
\begin{equation}
Re = \frac{\rho v L}{\eta} \sim
  \frac{10^3 \cdot 10^{-6} \cdot 10^{-5}}{10^{-3}}
  \sim 10^{-4},
\label{eq:reynolds}
\end{equation}
where $\rho \sim 10^3$\,kg/m$^3$ (water density),
$v \sim 1$\,$\mu$m/s (cytoplasmic streaming velocity),
$L \sim 10$\,$\mu$m (cell diameter), and
$\eta \sim 10^{-3}$\,Pa$\cdot$s (cytoplasmic viscosity).

At $Re \ll 1$, inertial terms are negligible and the Navier--Stokes
equations reduce to the Stokes equations:
\begin{equation}
-\nabla p + \eta \nabla^2 \mathbf{v} + \mathbf{f} = 0,
\label{eq:stokes}
\end{equation}
together with the continuity equation
$\nabla \cdot \mathbf{v} = 0$ (incompressibility).

The Kirchhoff-derived Eq.~\eqref{eq:continuous-kcl}, with
$G(\mathbf{r}) = \eta^{-1}$ and $\varphi = p$, is exactly the
Stokes equation in the form $\nabla \cdot (\eta^{-1} \nabla p)
= Q$, which combined with $\mathbf{v} = -\eta^{-1} \nabla p$
(Darcy's law in the Stokes limit) gives Eq.~\eqref{eq:stokes}.
\end{proof}

\begin{corollary}[Ray Tracer as Stokes Solver]
\label{cor:ray-stokes}
The ray tracer, by computing circuit currents along its path
(Eq.~\eqref{eq:ray-current}), is implicitly evaluating the
Stokes flow field. The accumulated current $\Jray$ along a ray
path is proportional to the line integral of the pressure gradient
along that path, which determines the flow velocity.
\end{corollary}


% ═════════════════════════════════════════════════════════════════════
\section{Viscosity from Partition Lag}
\label{sec:viscosity}

\begin{theorem}[Partition Lag Determines Viscosity]
\label{thm:viscosity}
The effective viscosity at position $\mathbf{r}$ is determined by the
partition lag:
\begin{equation}
\eta(\mathbf{r}) = \kb T(\mathbf{r}) \cdot \taup(\mathbf{r}) \cdot
  \rho(\mathbf{r}),
\label{eq:viscosity}
\end{equation}
where $\taup(\mathbf{r})$ is the local partition lag
(Eq.~\eqref{eq:partition-lag}), $T(\mathbf{r})$ is the local
temperature, and $\rho(\mathbf{r})$ is the local density.
\end{theorem}

\begin{proof}
The viscosity $\eta$ of a fluid relates the shear stress to the
velocity gradient: $\sigma_{xy} = \eta \cdot \partial v_x/\partial y$.
From the molecular kinetic theory perspective, viscosity arises from
the transfer of momentum between fluid layers through molecular
collisions. The relaxation time for this transfer is the time between
effective collisions, which in the partition framework is the partition
lag $\taup$ --- the time for a molecule to transition between
categorical states (i.e., to ``respond'' to the shear).

By the Green--Kubo relation from non-equilibrium statistical
mechanics:
\begin{equation}
\eta = \frac{V}{\kb T} \int_0^\infty
  \langle \sigma_{xy}(0) \sigma_{xy}(t) \rangle \, dt,
\label{eq:green-kubo}
\end{equation}
where $V$ is the volume and the angle brackets denote an equilibrium
ensemble average. The stress autocorrelation function decays with a
characteristic time $\taup$:
\begin{equation}
\langle \sigma_{xy}(0) \sigma_{xy}(t) \rangle \approx
  \langle \sigma_{xy}^2 \rangle \exp(-t/\taup).
\end{equation}

Substituting:
\begin{equation}
\eta \approx \frac{V}{\kb T} \langle \sigma_{xy}^2 \rangle \taup.
\end{equation}

Using the equipartition result $\langle \sigma_{xy}^2 \rangle =
(\kb T)^2 \rho / V$ (where $\rho$ is the number density), we obtain:
\begin{equation}
\eta \approx \kb T \cdot \rho \cdot \taup,
\end{equation}
which is Eq.~\eqref{eq:viscosity}.
\end{proof}

\begin{corollary}[Viscosity from Chromatographic Retention]
\label{cor:viscosity-retention}
Regions with high chromatographic retention (high $\tau \cdot \dS$)
have high viscosity. The ray is ``slowed down'' more in viscous regions,
exactly as a chromatographic analyte is retained more on a strongly
interacting stationary phase. This provides a direct measurement of
local viscosity from the ray's retention profile.
\end{corollary}

\begin{remark}
Experimentally, cytoplasmic viscosity is heterogeneous, ranging from
$\sim\!1$\,cP (water-like) in dilute regions to $\sim\!100$\,cP in
crowded regions near the cytoskeleton \citep{luby1999}. The partition
lag varies correspondingly: from $\sim\!10^{-12}$\,s in dilute regions
to $\sim\!10^{-10}$\,s in crowded regions. The ray tracer's retention
profile maps this heterogeneity directly.
\end{remark}

% ── Figure 11 proposal ─────────────────────────────────────────────
% FIGURE 11: Left: viscosity map of a cell (color coded from blue=low
% to red=high). Right: the same cell's chromatographic retention map
% from the ray tracer. The two maps are proportional, validating the
% partition-lag/viscosity correspondence.


% ═════════════════════════════════════════════════════════════════════
%             PART VII: VALIDATION AND APPLICATIONS
% ═════════════════════════════════════════════════════════════════════

\section{Validation Framework}
\label{sec:validation}

We propose five independent experimental designs to validate the
theory.

\subsection{Experiment 1: Triple Observation Consistency}

\textbf{Objective:} Verify the proportionality $\muabs \propto
1/(\tau \cdot \dS) \propto G \cdot RT$ on synthetic cellular volumes.

\textbf{Protocol:}
\begin{enumerate}[label=(\roman*),leftmargin=*]
\item Generate synthetic 3D cellular volumes with known partition
  states $\Sigma(x,y,z)$ at each voxel. Use the BBBC039 dataset
  \citep{bbbc039} cell segmentations as morphological templates and
  assign partition coordinates based on organelle identity.
\item Run the triple-observation ray march
  (Algorithm~\ref{alg:ray-march}) through each volume.
\item At each ray step, record $\muabs$, $\tau \cdot \dS$, and
  $G \cdot RT$ independently.
\item Compute the correlation coefficients between the three
  quantities across all positions and all volumes.
\end{enumerate}

\textbf{Expected result:} Pearson correlation $r > 0.95$ between all
three pairs, with the proportionality constants $\kappa_1$ and
$\kappa_2$ consistent with the theoretical predictions
(Eqs.~\eqref{eq:kappa1}--\eqref{eq:kappa2}).

\subsection{Experiment 2: Holographic Reconstruction Fidelity}

\textbf{Objective:} Verify that angular spectrum back-propagation from
2D to 3D recovers the ground-truth volumetric structure.

\textbf{Protocol:}
\begin{enumerate}[label=(\roman*),leftmargin=*]
\item Generate a synthetic 3D volume with known structure (e.g.,
  concentric spherical shells mimicking nuclear envelope, ER,
  mitochondria).
\item Forward-propagate to obtain the 2D holographic observation
  $U(x,y,0)$.
\item Back-propagate using Eq.~\eqref{eq:angular-spectrum} to
  reconstruct $U(x,y,z)$.
\item Extract partition coordinates using
  Eqs.~\eqref{eq:n-from-field}--\eqref{eq:s-from-field}.
\item Compare reconstructed partition coordinates with ground truth.
\end{enumerate}

\textbf{Expected result:} Reconstruction error $< 5\%$ in $n$ for
structures larger than $\lambda/2$ (the diffraction limit), with
degradation for sub-wavelength features.

\subsection{Experiment 3: Coherence Diagnostic Accuracy}

\textbf{Objective:} Verify that the interference visibility $\Vcell$
correctly distinguishes healthy and pathological cells.

\textbf{Protocol:}
\begin{enumerate}[label=(\roman*),leftmargin=*]
\item Generate synthetic volumes for healthy cells (synchronized
  oscillators: all eight classes at their nominal frequencies with
  small phase variance) and pathological cells (desynchronized:
  one or more classes with large phase variance or frequency shift).
\item Run the multi-ray interference pipeline
  (Algorithm~\ref{alg:multi-ray}) on each volume.
\item Compute $\Vcell$ and compare with the ground-truth coherence
  index (computed directly from the oscillator parameters).
\end{enumerate}

\textbf{Expected result:} $|\Vcell - \etacell| < 0.05$ for $N \geq 8$
ray directions. Healthy cells: $\etacell > 0.7$. Pathological cells:
$\etacell < 0.3$.

\subsection{Experiment 4: Flow Field Recovery}

\textbf{Objective:} Verify that the Doppler-shifted retention
(Eq.~\eqref{eq:doppler-retention}) correctly recovers a known
velocity field.

\textbf{Protocol:}
\begin{enumerate}[label=(\roman*),leftmargin=*]
\item Generate a synthetic volume with a known velocity field (e.g.,
  Poiseuille flow in a cylindrical domain, or a rotational flow
  mimicking cytoplasmic streaming).
\item Run rays in multiple directions through the moving medium.
\item Extract velocity using Eq.~\eqref{eq:flow-from-retention}.
\item Compare with the ground-truth velocity field.
\end{enumerate}

\textbf{Expected result:} Velocity recovery error $< 10\%$ for
$|\mathbf{v}|/\cS > 10^{-3}$ (i.e., velocities above
$\sim\!1$\,$\mu$m/s for $\cS \sim 10^3$\,m/s).

\subsection{Experiment 5: Real-Time Throughput}

\textbf{Objective:} Verify that the complete five-pass pipeline runs
at $\geq 20$ FPS on an integrated GPU.

\textbf{Protocol:}
\begin{enumerate}[label=(\roman*),leftmargin=*]
\item Implement the pipeline using WebGL~2.0
  \citep{khronos2024} or WebGPU \citep{webgpu2024}.
\item Run on three integrated GPUs: Intel UHD 630 ($\sim\!0.4$\,TFLOPS),
  AMD Radeon Vega 8 ($\sim\!1.1$\,TFLOPS), Apple M1 GPU
  ($\sim\!2.6$\,TFLOPS).
\item Measure per-pass and total latency. Report FPS.
\end{enumerate}

\textbf{Expected result:} $\geq 10$ FPS on Intel UHD 630, $\geq 20$
FPS on AMD Vega 8, $\geq 40$ FPS on Apple M1.

% ── Figure 12 proposal ─────────────────────────────────────────────
% FIGURE 12: Validation results summary. Five panels, one per
% experiment. (a) Scatter plot of mu_a vs 1/(tau*d_S) vs G*RT.
% (b) Reconstructed volume cross-sections. (c) ROC curve for
% coherence diagnostic. (d) Recovered vs ground-truth velocity field.
% (e) FPS bar chart across three integrated GPUs.


% ═════════════════════════════════════════════════════════════════════
\section{Discussion}
\label{sec:discussion}

\subsection{Convergence of Multiple Frameworks}

The ray-tracing framework presented here represents the convergence
of multiple independent theoretical developments:

\begin{enumerate}[leftmargin=*,itemsep=2pt]
\item \textbf{Oscillatory necessity and partition coordinates}
  provide the foundational structure: every persistent cellular process
  oscillates, and the state space is organized by partition coordinates
  $(n,\ell,m,s)$. These are the axioms from which everything follows
  (\S\ref{sec:axioms}--\S\ref{sec:partition}).

\item \textbf{Oxygen categorical microscopy} provides the ``invisible
  pixel'' --- the phase information that, combined with the visible
  amplitude pixel, creates the holographic record needed for 3D
  reconstruction (\S\ref{sec:dual-pixel}).

\item \textbf{Path enumeration and ray tracing} are mathematically
  identical operations. Enumerating all optical paths through a medium
  (as in Feynman's path integral formulation \citep{sakurai2017}) is
  the same as tracing all rays through a volume. The refraction puzzle
  --- how to reconstruct the 3D refractive index distribution from 2D
  observations --- is solved by holographic back-propagation
  (\S\ref{sec:angular-spectrum}).

\item \textbf{The dual-pixel stereogram} --- amplitude from one
  channel, phase from another --- is recognized as a hologram:
  it provides the complete wavefield needed for angular spectrum
  propagation (\S\ref{sec:dual-pixel}).

\item \textbf{Harmonic matching circuits} --- the idea that
  interference between oscillatory signals yields a diagnostic ---
  becomes the multi-ray interference calculation that produces the
  coherence index (\S\ref{sec:coherence}).

\item \textbf{S-distance chromatography} --- the cytoplasm as a
  column --- becomes the chromatographic leg of the triple observation,
  with retention encoding molecular interaction history
  (\S\ref{sec:chromatographic-transport}).

\item \textbf{Cellular observation equations} --- bioreactions as
  oscillations --- become the eight oscillator classes that determine
  ray--matter interactions
  (\S\ref{sec:oscillator-classes}).

\item \textbf{The cellular circuit model} --- reactions as circuit
  elements --- becomes the circuit leg of the triple observation,
  with Kirchhoff's laws extending to the Stokes equations for fluid
  flow (\S\ref{sec:circuit-transport}, \S\ref{sec:navier-stokes}).

\item \textbf{Universal spectral matching} --- all comparison as
  computer vision --- is realized by the GPU pipeline that performs
  the comparison operation as a rendering pass. The triple observation
  reduces molecular comparison to a visual comparison of ray-traced
  outputs.

\item \textbf{The GPU observation architecture} --- the fragment
  shader as measurement apparatus --- provides the computational
  substrate. The rendering--measurement identity means that the ray
  march is not a simulation of cellular observation; it \emph{is}
  the observation (\S\ref{sec:ray-march}).
\end{enumerate}

Ray tracing unifies all ten of these into a single volumetric
traversal: one ray march step simultaneously computes the result of
all ten frameworks.

\subsection{The Cellular Digital Twin}

The pipeline described in \S\ref{sec:pipeline} constitutes a
\emph{cellular digital twin} running in real time in a web browser.
From a single microscopy image:

\begin{enumerate}[leftmargin=*,itemsep=2pt]
\item The holographic back-propagation reconstructs the 3D
  morphology.
\item The optical observation produces a volumetric absorption map
  (identifying chromophore-rich regions: mitochondria, ER, nucleus).
\item The chromatographic observation produces a metabolic flow map
  (identifying regions of high molecular interaction: active enzyme
  complexes, transport vesicles).
\item The circuit observation produces a reaction network state map
  (identifying metabolic hotspots and bottlenecks).
\item The coherence index provides a single scalar health diagnostic.
\item The flow field recovery provides cytoplasmic velocity maps.
\end{enumerate}

All of this from a single 2D image, in $\sim\!43$\,ms, on an
integrated GPU, in a browser.

\subsection{Implications for Personalized Medicine}

The cellular digital twin has direct implications for personalized
medicine. A patient's cells, imaged by standard fluorescence microscopy,
can be analyzed in real time by the ray-tracing pipeline. The coherence
index $\etacell$ provides an immediate, quantitative health diagnostic
without requiring expert pathologist review. Specific decoherent
oscillator classes identify the type of pathology (e.g., mitochondrial
dysfunction from ATP oscillator decoherence, calcium signaling disorder
from Ca oscillator decoherence). The metabolic flow map identifies
drug targets (reaction bottlenecks). The 3D reconstruction enables
virtual biopsy.

The computation requires no cloud infrastructure, no discrete GPU, and
no specialized software beyond a web browser. This makes the technology
accessible to any clinic with a microscope and a laptop.

\subsection{Coupled Oscillator Perspective}

The eight oscillator classes do not operate independently within a
living cell. They are coupled through shared molecular substrates:
calcium oscillations modulate enzyme activity, which modulates ATP
production, which modulates membrane potential, which modulates channel
gating. This coupling is precisely the type described by the Kuramoto
model of coupled oscillators \citep{kuramoto1975}, extended to a
heterogeneous population spanning 18 decades of frequency.

The coherence index $\etacell$ measures the \emph{synchronization
order parameter} of this coupled oscillator system. A healthy cell has
a high order parameter (oscillators synchronized despite diverse
frequencies), analogous to a Kuramoto system above the critical
coupling threshold. Disease corresponds to a drop below the
synchronization threshold, detectable as a decrease in ray
interference visibility.

\subsection{Autopoietic Organization}

The self-maintaining character of living cells --- their autopoiesis
\citep{varelamaturana1980} --- manifests in the ray-tracing framework
as a self-consistent volumetric partition state. The 3D partition field
$\Sigma(\mathbf{r})$ determines the ray--matter interactions, which in
turn determine the metabolic fluxes, which in turn maintain the
partition field. This circular causality is a defining feature of
autopoietic systems and is captured naturally by the iterative nature
of the ray march: each step updates the ray state based on the medium,
and the medium state is maintained by the cumulative effect of all
molecular interactions (currents) the ray computes.

\subsection{Autocatalytic Sets and Network Self-Organization}

The circuit model of \S\ref{sec:circuit-transport} maps directly onto
the theory of autocatalytic sets \citep{kauffman1986}. An autocatalytic
set is a collection of molecular species where each species'
production is catalyzed by at least one other species in the set.
In the circuit model, this corresponds to a strongly connected
component of the reaction graph where every node has at least one
incoming edge with positive current. The ray tracer, by computing
currents along its path, can identify autocatalytic sets as connected
regions where $\Jray$ is consistently positive in all directions.

\subsection{Comparison with Stochastic Simulation}

The Gillespie algorithm \citep{gillespie1977} provides exact
stochastic simulation of chemical reaction networks. However, it
operates on a well-mixed assumption and scales as $O(N_{\mathrm{rxn}}
\cdot N_{\mathrm{steps}})$ per trajectory. The ray-tracing approach
is complementary: it provides a spatially resolved, deterministic
(mean-field) solution in $O(V^3 \cdot N_{\mathrm{march}})$ time,
where $V^3$ is the volume resolution and $N_{\mathrm{march}}$ is the
number of ray march steps. For a $128^3$ volume with 256 steps, this
is $\sim\!5 \times 10^8$ operations --- comparable to a single
Gillespie trajectory of $10^6$ reactions for a network of $\sim\!500$
species. The ray tracer additionally provides the 3D spatial
information that Gillespie lacks.

\subsection{Limitations}

Several limitations should be acknowledged:

\begin{enumerate}[leftmargin=*,itemsep=2pt]
\item \textbf{Resolution of 3D reconstruction.} The angular spectrum
  back-propagation is diffraction-limited: structures smaller than
  $\sim\!\lambda/2$ cannot be resolved. For visible light ($\lambda
  \approx 500$\,nm), this limits resolution to $\sim\!250$\,nm.
  Sub-cellular structures smaller than this (e.g., individual proteins
  at $\sim\!5$\,nm) are not resolved but contribute to the averaged
  partition state of each voxel.

\item \textbf{Model assumptions.} The triple observation identity
  (Theorem~\ref{thm:triple-observation}) is derived under a
  leading-order approximation where $n/n_{\max} \ll 1$. For highly
  concentrated regions where $n \to n_{\max}$, higher-order
  corrections may be needed. The linear chromatographic model
  (retention proportional to S-distance) may break down at very high
  analyte concentrations (column overloading).

\item \textbf{Phase retrieval.} The invisible pixel (O$_2$ emission
  timing) provides phase information, but the signal-to-noise ratio
  may be limited in practice. Phase retrieval algorithms (e.g.,
  Gerchberg--Saxton iteration) can supplement direct phase measurement
  but add computational cost.

\item \textbf{Temporal resolution.} The pipeline runs at
  $\sim\!23$\,FPS, which resolves dynamics down to $\sim\!40$\,ms.
  This captures circadian, genetic, calcium, and ATP oscillators but
  not the faster oscillator classes (channel, enzyme, protein). These
  faster classes are captured only through their time-averaged effect
  on the partition state.

\item \textbf{Multi-cell extension.} The current framework treats a
  single cell. Extension to tissue-scale (multi-cell) ray tracing
  requires either larger volumes (higher memory) or hierarchical
  approaches (coarse outer volume, fine per-cell volumes).
\end{enumerate}

\subsection{Future Directions}

\begin{enumerate}[leftmargin=*,itemsep=2pt]
\item \textbf{Learned ray interaction functions.} The form factor
  $F(\ell,m,s)$ and the S-distance weights $\mathbf{w}$ can be
  treated as learnable parameters, optimized by GPU-supervised
  training using the coherence index and other self-supervised
  signals.

\item \textbf{Multi-cell ray tracing.} Extend to tissue-scale
  computation by treating each cell as a node in a macro-circuit,
  with intercellular signaling (gap junctions, paracrine signals)
  as inter-node conductances.

\item \textbf{Temporal ray tracing.} Replace the static volume with a
  time-varying volume, updated each frame from new microscopy data.
  The ray tracer then tracks the evolution of the cellular state over
  time, enabling dynamic analysis.

\item \textbf{Viterbi decoding of reaction history.} The phase
  accumulation along a ray path encodes a sequence of partition
  transitions. This sequence can be decoded using the Viterbi
  algorithm \citep{viterbi1967} to find the most likely sequence of
  bioreactions along the path, providing a ``molecular movie'' from a
  single static image.

\item \textbf{WebGPU implementation.} The current design targets
  WebGL~2.0 for maximum compatibility. WebGPU \citep{webgpu2024}
  provides compute shaders, which would enable more efficient FFT
  (for holographic back-propagation) and parallel reduction (for
  scalar readback).
\end{enumerate}


% ═════════════════════════════════════════════════════════════════════
\section{Conclusion}
\label{sec:conclusion}

We have proved from first principles that a single ray marching through
a cellular volume simultaneously computes three physically distinct
observations --- optical absorption, chromatographic retention, and
biochemical circuit current --- at every integration step. The Triple
Observation Identity (Theorem~\ref{thm:triple-observation}) establishes
that these three are not independent computations but the same
partition observation expressed in three equivalent representations.
The proportionality $\muabs \propto 1/(\tau \cdot \dS) \propto G \cdot
RT$ is a consequence of the Triple Equivalence Theorem
(\S\ref{sec:triple-equivalence}), which proves that oscillatory,
categorical, and partitional descriptions are mathematically identical.

The eight cellular oscillator classes spanning 18 decades of frequency
--- from protein folding vibrations at $10^{13}$\,Hz to circadian
rhythm at $10^{-5}$\,Hz --- each manifest as a specific type of
ray--matter interaction: absorption, scattering, gated transmission,
reflection, amplification, state switching, oscillatory modulation,
and parametric drift. The accumulated phase shift along a ray path
encodes the complete reaction history. Interference between multiple
rays yields a visibility parameter that is identical to the cellular
coherence index: a single scalar diagnostic of cellular health.

The holographic reconstruction from a 2D observation, via angular
spectrum back-propagation, provides the 3D volumetric partition state
at every spatial position. The five-pass GPU pipeline --- dual-pixel
analysis, holographic back-propagation, triple-observation ray march,
multi-ray interference, and scalar readback --- runs at
$\sim\!23$ frames per second on an integrated GPU with
$\sim\!45$\,MB working set.

The ray tracer is not a simulation of the cell. By the
rendering--measurement identity, the act of ray tracing \emph{is} the
act of observation. The rendered output is the measurement result
itself: a cellular digital twin running in a web browser.

The extension to cytoplasmic fluid dynamics, via the Kirchhoff--Stokes
correspondence, shows that the ray tracer implicitly solves the
low-Reynolds-number flow equations from the circuit currents it
computes. Local viscosity is measured from chromatographic retention.
Flow velocity is measured from retention anisotropy.

The framework unifies optical microscopy, chromatographic analysis,
systems biology, fluid dynamics, and diagnostic medicine into a single
computational operation: the volumetric ray march. This convergence is
not accidental; it is a mathematical consequence of the Triple
Equivalence, which proves that nature has one underlying structure
(the partition state) that we observe through different instruments
(microscope, chromatograph, multimeter) and mistakenly believe to be
different phenomena. The ray tracer computes all simultaneously because
they are the same thing.


% ═════════════════════════════════════════════════════════════════════
%                         REFERENCES
% ═════════════════════════════════════════════════════════════════════

\bibliographystyle{plainnat}
\bibliography{references}

\end{document}
